\documentclass[12pt]{article}

\usepackage[a4paper, margin=2.5cm]{geometry}
\usepackage{amsmath}
\usepackage{amssymb}
\usepackage{amsthm}
\usepackage[colorlinks=true, linkcolor=blue, citecolor=blue, urlcolor=blue]{hyperref}
\usepackage{booktabs}
\usepackage{natbib}

\newtheorem{theorem}{Theorem}
\newtheorem{proposition}[theorem]{Proposition}
\newtheorem{corollary}[theorem]{Corollary}
\newtheorem{definition}[theorem]{Definition}
\newtheorem{remark}[theorem]{Remark}
\newtheorem{example}[theorem]{Example}
\newtheorem{lemma}[theorem]{Lemma}

\newcommand{\Ecal}{\mathcal{E}}
\newcommand{\Gcal}{\mathcal{G}}
\newcommand{\Fcal}{\mathcal{F}}
\newcommand{\Ocal}{\mathcal{O}}
\newcommand{\Ccal}{\mathcal{C}}
\newcommand{\Hcal}{\mathcal{H}}
\newcommand{\Sadm}{\mathcal{S}_{\mathrm{adm}}}
\newcommand{\Mcl}{\mathcal{M}_{\mathrm{cl}}}
\newcommand{\Rst}{R_{\mathrm{st}}}
\newcommand{\Pst}{P_{\mathrm{st}}}

\title{Probability as Constraint Geometry:\\[6pt]
\large Toward a Derivation of the Born Rule from Closure Dynamics}
\author{%
  Lee Smart\\[4pt]
  \textit{Institute of Vibrational Field Dynamics}\\[2pt]
  \texttt{contact@vibrationalfielddynamics.org}\\[2pt]
  \texttt{@vfd\_org}%
}
\date{April 2026}

\begin{document}

\maketitle

\noindent\textbf{Status:} Preprint (not peer-reviewed)

\begin{abstract}
Paper~XXIX defined the observer as a constraint substructure within the VFD event-order framework and showed that measurement corresponds to constraint-induced reduction of admissible configurations. The present paper addresses the next gap: \textit{why do admissible outcomes have different probabilities, and why do those probabilities take the quadratic form $P_i = |\psi_i|^2$?}

Admissibility alone determines which outcomes are possible; it does not determine their relative frequencies. This paper proposes that outcome probabilities arise as geometric weights over the set of observer-admissible configurations, with weights tied to closure dynamics only after additional measurement and observer hypotheses (the sampling postulate and state-completeness of Route~C; universal completeness and basis-realisability of Route~E) are imposed.

Five weighting routes are examined. Routes~A and~B (basin volume, stability depth) are motivational prototypes illustrating why probability must be geometric. Route~C --- the paper's central conditional derivation --- shows that in the equilibrium tangent regime of Papers~XVIII--XXI, the Nelson density $|\Psi|^2$ coincides with the Paper~XIV stationary distribution $\Pst \propto e^{-2\Fcal/\sigma^2}$, so that sector probabilities reduce to $P_i = |c_i|^2$ only under five stated assumptions (equilibrium, well-separated sectors, sector-supported decomposition, the sampling hypothesis that measurement frequencies equal integrated stationary sector masses, and state-completeness of the observer for the state in question). Without state-completeness, the result is the observer-conditional ratio $|c_i|^2/\sum_{j\in\mathrm{adm}}|c_j|^2$. Route~D confirms uniqueness via Gleason's theorem once Hilbert structure is granted. Route~E --- the central new contribution --- packages imported equilibrium-tangent geometric inputs (the $U(N)$-invariant K\"ahler structure on $\mathbb{CP}^{N-1}$ supplied by Paper~XXI under its sector-subspace invariance hypothesis; the Nelson-wavefunction complex-structure bookkeeping of Paper~XVIII; basin locality and LMC from Paper~XXIX, upgraded here to a restricted, state-specific non-contextuality theorem under completeness) together with a local Laplace identification of the closure Hessian with the Fisher information metric of a location-parameter family on configuration space, and applies Petz's uniqueness and a Gleason-type frame-function argument. Theorem~\ref{thm:born_unique} itself is a Gleason theorem: under axioms (U1)--(U4) on a global probability assignment $P$ on $\mathbb{CP}^{N-1}$ --- full $U(N)$-covariance of $P$, the global frame-function condition, continuity, and eigenstate certainty --- together with universal completeness $\Mcl(\Ocal) = \mathbb{CP}^{N-1}$, the Born assignment is singled out as the unique compatible probability for $N \geq 3$ (the $N = 2$ case is not covered, since Gleason's theorem requires $\dim \geq 3$; see Remark~\ref{rem:N_two}). Going from VFD's observer-indexed conditional probabilities to such a single global $P$ requires an additional bridge: the basis-realisability assumption that every orthonormal basis is realised by a complete localised-probe observer (Remark~\ref{rem:basis_realisability}), which, combined with universal completeness and Paper~XXIX's LMC-derived restricted non-contextuality, yields the global (U2). Paper~XXI supplies the imported equilibrium-tangent unitary/projective geometric action; covariance of $P$ under that action is the (U1) hypothesis on $P$ itself, assumed and not derived. Route~E is therefore a \emph{conditional uniqueness argument} given imported structures, not a derivation of the Born rule from closure alone.

The result is a \emph{conditional reconstruction} of Born weighting under imported geometric structure and additional observer-side hypotheses. Given (i) Paper~XXI supplying the $U(N)$-invariant K\"ahler structure on $\mathbb{CP}^{N-1}$ as standard geometric QM on an assumed invariant subspace, (ii) Paper~XVIII supplying the Nelson-wavefunction complex-structure bookkeeping, (iii) Paper~XXIX supplying the localised-probe-observer variant and its local measurement consistency theorem, (iv) the Laplace regime supplying a candidate local scale via Proposition~\ref{prop:hessian_fisher} (its use as a global scale-fix for Fubini--Study is conditional on a bridge theorem not proved here), (v) universal completeness $\Mcl(\Ocal) = \mathbb{CP}^{N-1}$, and (vi) axioms (U1)--(U4) of Theorem~\ref{thm:born_unique} --- Paper~XXI supplies the equilibrium-tangent unitary/projective geometric action, with covariance of $P$ assumed as hypothesis (U1), the global frame-function hypothesis (U2) taken as a hypothesis, continuity, and eigenstate certainty --- the Born measure is the unique probability assignment compatible with the construction for $N \geq 3$. A restricted, state-specific form of non-contextuality is derivable within Paper~XXIX's localised-probe-observer variant from basin locality and LMC; upgrading to the global (U2) requires, in addition, universal completeness and the assumption that every orthonormal basis is realised by such an observer.

A structural feature of the VFD framework worth highlighting: the probability assignment is intrinsically \emph{conditional}, normalised over the observer's admissible configuration mass. For \emph{complete} observers --- those whose admissible set covers the state's Nelson density --- this reduces to the unconditional standard Born rule. For \emph{partial} observers --- whose admissible set omits some of the state's support --- VFD predicts context-dependent conditional probabilities. Under the sector-supported decomposition used in Route~C, these conditional probabilities take the explicit form $|c_i|^2 / \sum_{j \in \mathrm{adm}}|c_j|^2$, with the complement $1 - \sum_{j \in \mathrm{adm}}|c_j|^2$ interpreted as the mass of non-detection by this observer. This is a formal extension of the paper's observer model beyond the complete-measurement regime --- its physical scope remains open --- and not a contradiction of standard QM: for complete observers the two prescriptions coincide exactly.
\end{abstract}

%==================================================================
\section{Introduction}\label{sec:intro}
%==================================================================

The Born rule --- the prescription that the probability of measuring outcome $i$ is $P_i = |\langle \phi_i | \psi \rangle|^2$ --- is one of the most empirically successful postulates in physics. It is also one of the least explained. Standard quantum mechanics introduces it as an axiom, without providing a mechanism for why measurement outcomes should be weighted by the squared modulus of the amplitude.

\subsection*{The gap}

The VFD programme has addressed several foundational questions in sequence:
\begin{itemize}
    \item Papers~XII--XIV: closure dynamics, stochastic extension, stationary distribution.
    \item Papers~XV--XXI: quantum structure recovery (Schr\"odinger as equilibrium tangent limit).
    \item Paper~XXII: spectral bridge connecting mass, gauge, and constants.
    \item Papers~XXIII--XXVII: gravitational structure from event-order geometry.
    \item Paper~XXVIII: unification bridge between quantum and gravitational tracks.
    \item Paper~XXIX: observer placement as constraint substructure.
\end{itemize}

Paper~XXIX showed that measurement corresponds to constraint-induced reduction of admissible configurations. But admissibility is binary: a configuration is either admissible or not. This does not explain why different admissible outcomes are realised with different frequencies. The Born-rule gap remains:

\begin{quote}
Given a set of observer-admissible outcomes, what determines their relative probabilities?
\end{quote}

\subsection*{Proposal}

This paper proposes that outcome probabilities are not primitive stochastic ingredients added to measurement. They are the induced measure over observer-compatible admissible configurations, with weights determined by the geometry and dynamics of constraint satisfaction.

\subsection*{Goals}

This paper has two goals:
\begin{enumerate}
    \item Explain why admissibility must be supplemented by a weighting measure to account for observed frequencies.
    \item Identify the equilibrium stationary measure of the closure dynamics as the natural candidate for that weighting.
\end{enumerate}

\subsection*{Scope}

The paper develops five routes toward quadratic weighting, works a toy model explicitly, and identifies the conditions under which the Born rule emerges. Routes~A and~B are motivational prototypes. Route~C is a conditional dynamical derivation: in the equilibrium tangent regime, sector probabilities are inherited from the stationary measure. Route~D is Gleason-style uniqueness given Hilbert structure. Route~E packages imported geometric inputs (the Paper~XVIII complex-structure bookkeeping, the Paper~XXI equilibrium-tangent $U(N)$-invariant K\"ahler structure on $\mathbb{CP}^{N-1}$, and a local Laplace-regime closure-Hessian/Fisher computation) and applies Gleason's frame-function theorem. The uniqueness proof itself is Gleason's; the role of the Chentsov--Petz content is structural (identifying the imported K\"ahler structure as Fubini--Study up to scale) rather than load-bearing inside the uniqueness proof. A restricted, state-specific form of non-contextuality is derived here from Paper~XXIX's basin-locality and LMC results within the complete localised-probe-observer variant; global (U2) as used by Theorem~\ref{thm:born_unique} remains an added hypothesis unless universal completeness and basis-realisability are also imposed (Remark~\ref{rem:basis_realisability}). The paper is explicit about which inputs are imported from prior papers.

The paper distinguishes three levels of claim:
\begin{itemize}
    \item[\textbf{Level I.}] Results from prior VFD papers taken as given: the stochastic closure dynamics and its stationary distribution (Paper~XIV), the Nelson pairing and Schr\"odinger equilibrium (Papers~XVIII and~XXI), and the observer/admissibility framework (Paper~XXIX).
    \item[\textbf{Level II.}] The conditional derivation in this paper (Route~C): in the equilibrium tangent regime, with well-separated sectors, sector-supported decomposition, the sampling hypothesis (measurement frequencies equal integrated stationary mass over sectors; Proposition~\ref{prop:born} assumption (4)), and completeness of the observer for the state, $P_i = |c_i|^2$ holds up to overlap corrections.
    \item[\textbf{Level III.}] The uniqueness programme in this paper (Route~E): given Levels~I and~II, the Laplace regime, the stipulation that physical measurements realise Paper~XXIX's localised-probe-observer variant, and the additional hypotheses of Theorem~\ref{thm:born_unique} --- full $U(N)$-covariance, global (U2), continuity, eigenstate certainty, universal completeness $\Mcl(\Ocal) = \mathbb{CP}^{N-1}$, and basis-realisability --- the Born measure is uniquely selected. Non-contextuality of probability is derived only in a restricted, state-specific form from Paper~XXIX's LMC result (Section~\ref{subsec:noncontextuality}); upgrading to the global frame-function hypothesis (U2) requires both universal completeness and basis-realisability (Remark~\ref{rem:basis_realisability}).
\end{itemize}
Levels~I--III are stated separately so that the conditional scope of the result is transparent.

%==================================================================
\section{Programme Position}\label{sec:position}
%==================================================================

The logical chain from prior papers:

\begin{enumerate}
    \item \textbf{Structure.} The 600-cell provides the geometric substrate (Papers~I--V, XXII).
    \item \textbf{Dynamics.} Closure gradient flow and stochastic extension (Papers~XII--XIV).
    \item \textbf{Quantum recovery.} Schr\"odinger equation as equilibrium tangent limit (Papers~XV--XXI).
    \item \textbf{Observer.} Bounded, stable, self-referential constraint substructure (Paper~XXIX).
    \item \textbf{Measurement.} Reduction to observer-admissible configurations (Paper~XXIX).
    \item \textbf{Probability.} Geometric weighting over admissible sectors (\textit{this paper}).
\end{enumerate}

Step 6 is the natural continuation: once measurement is understood as admissibility reduction, the remaining task is to explain the weighting.

%==================================================================
\section{Formal Setup}\label{sec:setup}
%==================================================================

\subsection{Spaces and notation}\label{subsec:notation}

Several distinct spaces appear throughout the paper. We fix notation for each now and use it consistently:

\begin{center}
\renewcommand{\arraystretch}{1.15}
\begin{tabular}{@{}p{0.25\textwidth}p{0.3\textwidth}p{0.4\textwidth}@{}}
\toprule
Name & Symbol & What it contains \\
\midrule
Configuration space & $\mathbb{R}^\Ecal$ & Field configurations $\Phi: \Ecal \to \mathbb{R}$ on the event set $\Ecal$. An element is denoted $\Phi \in \mathbb{R}^\Ecal$. \\
Admissible set & $\Phi_{\mathrm{adm}}(\Ocal) \subset \mathbb{R}^\Ecal$ & Configurations passing the observer coherence constraint (Definition~\ref{def:admissible}). \\
Closure basin & $B_{\phi_i} \subset \mathbb{R}^\Ecal$ & Basin of attraction under the gradient flow of $\Fcal$ converging to the minimum $\phi_i$ (Paper~XXIX). \\
Sector-superposition Hilbert space & $\Hcal_\Ocal \cong \mathbb{C}^N$ & Finite-dimensional complex Hilbert space of sector superpositions $\Psi = \sum_i c_i \phi_i$ (Paper~XXI). \\
Projective state manifold & $\mathbb{CP}^{N-1}$ & Pure states of $\Hcal_\Ocal$ modulo global phase. $[\Psi] \sim [e^{i\alpha}\Psi]$. \\
Closure statistical parameter space & $\Mcl(\Ocal) \subseteq \mathbb{CP}^{N-1}$ & The admissible subset of pure projective states (Definition~\ref{def:closure_manifold}); treated formally as a statistical parameter space, not a proved smooth manifold. \\
\bottomrule
\end{tabular}
\end{center}

The key structural link between the Hilbert-space side and the configuration-space side is the Nelson pairing of Paper~XVIII: a state $\Psi \in \Hcal_\Ocal$ defines a density $\rho_\Psi = |\Psi|^2$ on $\mathbb{R}^\Ecal$. Reading the sector superposition $\Psi = \sum_i c_i \phi_i$ as a function on configuration space with sector-localised components $\phi_i$ is a \emph{modelling assumption} of this paper (see Remark following Definition~\ref{def:closure_manifold}); Paper~XVIII does not itself construct such a sector-localised basis. Sector probabilities in this paper are the \emph{normalised conditional masses}
\begin{equation}\label{eq:prob_via_integration}
    P(\Pi_i, \Psi; \Ocal) \;=\; \frac{\int_{B_{\phi_i} \cap \Phi_{\mathrm{adm}}(\Ocal)} |\Psi(\Phi)|^2\, d\mu(\Phi)}{\int_{\Phi_{\mathrm{adm}}(\Ocal)} |\Psi(\Phi)|^2\, d\mu(\Phi)},
\end{equation}
where $\Pi_i = |\phi_i\rangle\langle\phi_i|$ is the rank-1 projector onto the sector minimum and the observer $\Ocal$ enters through its admissibility set $\Phi_{\mathrm{adm}}(\Ocal)$. Physically, $P(\Pi_i, \Psi; \Ocal)$ is the conditional probability of outcome $\Pi_i$ \emph{given that this observer registers some outcome}; the denominator normalises over the observer-admissible configuration mass.

Two regimes deserve separation. Fix a state $\Psi \in \Hcal_\Ocal$. We say the observer $\Ocal$ is \emph{complete for $\Psi$} if the Nelson density $\rho_\Psi$ is supported in $\Phi_{\mathrm{adm}}(\Ocal)$, equivalently if the denominator of~\eqref{eq:prob_via_integration} equals $\|\Psi\|^2 = 1$. We say $\Ocal$ is \emph{(universally) complete} if it is complete for every pure state, equivalently if $\Mcl(\Ocal) = \mathbb{CP}^{N-1}$. In either complete regime for $\Psi$, $P$ reduces to the unconditional sector mass; universal completeness is the stronger hypothesis that Theorem~\ref{thm:born_unique} imposes. Standard quantum mechanics implicitly works in the completeness regime: a measurement basis covers the state, Parseval gives $\sum_i |c_i|^2 = 1$, and normalised conditional probability equals the raw sector mass. When $\Ocal$ is \emph{partial for $\Psi$} --- admitting only some basins while $\rho_\Psi$ has weight elsewhere --- the denominator of~\eqref{eq:prob_via_integration} is strictly less than $\|\Psi\|^2$, and VFD's conditional probability differs from the unconditional Born weight. The partial case is a formal extension of the present observer formalism beyond the complete-measurement regime; its physical scope remains open (see Remark~\ref{rem:complete_partial} below).

\subsection{Configuration space and admissibility}

The closure functional $\Fcal: \mathbb{R}^\Ecal \to \mathbb{R}_{\geq 0}$ assigns a closure cost to each configuration.

\begin{definition}[Observer-admissible configurations, restated from Paper~XXIX]\label{def:admissible}
Given an observer $(\Ocal, \Ccal_\Ocal)$ as defined in Paper~XXIX, the admissible set is
\begin{equation}\label{eq:admissible}
    \Phi_{\mathrm{adm}}(\Ocal) = \bigl\{\, \Phi \in \mathbb{R}^\Ecal \;\big|\; \Ccal_\Ocal\!\bigl(\Phi|_\Ocal, \partial_\Ocal \Phi\bigr) \leq \epsilon \,\bigr\},
\end{equation}
where $\epsilon > 0$ is a coherence threshold. This matches Paper~XXIX's admissibility definition. Note that Paper~XXIX additionally requires an $\Fcal$-local-minimum condition for the combined measurement condition ($\Fcal$ locally minimal \emph{and} $\Ccal_\Ocal \leq \epsilon$); in the present paper we work with coherence-threshold admissibility alone and impose local-minimum conditions separately where needed.
\end{definition}

\subsection{Outcome sectors}

When the admissible set contains multiple locally stable configurations, we partition it into \emph{outcome sectors} corresponding to basins of attraction. With a slight reuse of the symbol from the dictionary (the subscripted version $\Phi_i$ is always the sector subset of configuration space, never an element):
\begin{equation}\label{eq:sector_partition}
    \Phi_{\mathrm{adm}}(\Ocal) \;=\; \bigsqcup_{i=1}^{N} \Phi_i(\Ocal), \qquad \Phi_i(\Ocal) \;:=\; B_{\phi_i} \cap \Phi_{\mathrm{adm}}(\Ocal),
\end{equation}
where $B_{\phi_i}$ is the basin of attraction of the $i$-th local minimum $\phi_i$ of $\Fcal$. The sector subset $\Phi_i(\Ocal) \subset \mathbb{R}^\Ecal$ is \emph{not} the same object as the sector basis vector $\phi_i \in \Hcal_\Ocal$ or the sector projector $\Pi_i = |\phi_i\rangle\langle\phi_i|$; the three are related but typed differently:
\begin{equation}
    \underbrace{\Phi_i(\Ocal)}_{\subset\,\mathbb{R}^\Ecal}
    \quad\longleftrightarrow\quad
    \underbrace{\phi_i}_{\in\,\Hcal_\Ocal}
    \quad\longleftrightarrow\quad
    \underbrace{\Pi_i = |\phi_i\rangle\langle\phi_i|}_{\in\,\mathrm{End}(\Hcal_\Ocal)}.
\end{equation}
Paper~XVIII's Nelson pairing provides the relation between Hilbert-space states $\Psi$ and configuration-space densities $\rho_\Psi = |\Psi|^2$. In the sector-supported decomposition regime assumed by Route~C, the sector basis vectors $\phi_i \in \Hcal_\Ocal$ are chosen as normalised modes whose Nelson densities $|\phi_i|^2$ are supported (up to exponentially suppressed tails) on the basins $\Phi_i(\Ocal)$; this choice is a modelling assumption consistent with, but not a specific construction provided by, Paper~XVIII. Equation~\eqref{eq:prob_via_integration} then reads the integrated density over sector $\Phi_i(\Ocal)$ as the probability associated with projector $\Pi_i$.

\subsection{The probability question}

Admissibility determines which sectors are non-empty. It does not determine their relative likelihood:
\begin{equation}\label{eq:not_equal}
    \Phi_i(\Ocal) \neq \emptyset \;\;\text{for all } i \quad \not\Rightarrow \quad P(\Pi_i, \Psi; \Ocal) = \tfrac{1}{N} \;\;\text{for all } i.
\end{equation}
The task is to construct a probability measure over outcome sectors from the structure of the closure dynamics.

\subsection{Complete and partial observers: a structural feature of the VFD framework}\label{subsec:complete_partial}

The probability assignment~\eqref{eq:prob_via_integration} is \emph{conditional} by construction: it normalises over the observer-admissible configuration mass. This splits the observer class into two regimes, which deserve explicit separation because they correspond to two physically distinct measurement situations.

\begin{remark}[Complete observers and recovery of standard Born]\label{rem:complete_partial}
An observer $\Ocal$ is \emph{complete for a state $\Psi$} if the Nelson density $\rho_\Psi = |\Psi|^2$ is entirely supported within $\Phi_{\mathrm{adm}}(\Ocal)$. This is a state-specific condition. An observer is \emph{universally complete} if $\Mcl(\Ocal) = \mathbb{CP}^{N-1}$, i.e.\ it is complete for every pure state in $\Hcal_\Ocal$; universal completeness is the strictly stronger property. In either completeness regime for $\Psi$,
\begin{equation}\label{eq:complete_observer}
    \int_{\Phi_{\mathrm{adm}}(\Ocal)} |\Psi|^2\, d\mu \;=\; \|\Psi\|^2 \;=\; 1,
\end{equation}
so the denominator of~\eqref{eq:prob_via_integration} equals $1$ and $P(\Pi_i, \Psi; \Ocal)$ reduces to the unconditional sector mass $|c_i|^2$. This is the standard Born-rule regime: a measurement apparatus that covers every outcome the state can yield. Theorem~\ref{thm:born_unique} below imposes the universal-completeness hypothesis $\Mcl(\Ocal) = \mathbb{CP}^{N-1}$, which is strictly stronger than state-completeness for a single $\Psi$ and is load-bearing for Gleason's frame-function argument (Remark~\ref{rem:full_CP_hypothesis}).
\end{remark}

\begin{remark}[Partial observers and VFD's conditional-probability prediction]\label{rem:partial}
An observer is \emph{partial} for a state $\Psi$ if $\rho_\Psi$ has weight outside $\Phi_{\mathrm{adm}}(\Ocal)$ --- that is, there exist basins in the support of $\Psi$ that the observer does not distinguish or admit. In this regime the denominator of~\eqref{eq:prob_via_integration} is strictly less than $\|\Psi\|^2$, and
\begin{equation}\label{eq:partial_observer}
    P(\Pi_i, \Psi; \Ocal) \;=\; \frac{|c_i|^2}{\sum_{j \in \mathrm{adm}(\Ocal)} |c_j|^2} \;\neq\; |c_i|^2,
\end{equation}
where $\mathrm{adm}(\Ocal)$ indexes the basins admitted by $\Ocal$. The VFD prediction is context-dependent conditional probability: $P(\Pi_i, \Psi; \Ocal)$ is the probability of outcome $i$ given that a measurement by $\Ocal$ registered some admissible outcome. This is not the standard Born rule; it is the conditional distribution naturally associated with a partial measurement.
\end{remark}

\begin{remark}[Relation to standard quantum mechanics]\label{rem:vs_standard_qm}
Standard quantum mechanics assigns unconditional probabilities $|c_i|^2$ to outcomes of a measurement and implicitly assumes the measurement is complete --- every outcome the state can yield is covered by some projector of the PVM. Partial measurements in standard QM are typically handled by attaching a ``no-click'' outcome or by extending the PVM. VFD's observer formalism handles this differently: it assigns a conditional probability~\eqref{eq:partial_observer} given that the observer registers, and leaves the mass $1 - \sum_{j \in \mathrm{adm}}|c_j|^2$ of unadmitted support as a structural prediction of \emph{non-detection} for this specific observer.

This is not a contradiction with standard QM. For complete observers the two prescriptions coincide exactly. For partial observers VFD predicts a clean conditional-probability structure that standard QM either hides inside the POVM formalism or absorbs into ad hoc no-click additions. The cost of this clarity is that the theorem statements below must carry the completeness hypothesis $\Mcl(\Ocal) = \mathbb{CP}^{N-1}$ where they intend to reproduce the standard Born rule; the benefit is that the formalism naturally extends to partial measurements without extra machinery.

This generalisation is a formal feature of the present VFD observer formalism; its physical scope remains open.
\end{remark}

%==================================================================
\section{Why Admissibility Alone Is Not Enough}\label{sec:gap}
%==================================================================

Consider a system with two admissible outcome sectors $\Phi_1$ and $\Phi_2$. Both satisfy the observer constraint. Both are locally stable under the closure functional. If admissibility were the whole story, the prediction would be $P_1 = P_2 = 1/2$ for any two-outcome measurement.

This contradicts experiment. Quantum mechanics predicts --- and experiment confirms --- that outcome probabilities depend on the state: $P_i = |\langle \phi_i | \psi \rangle|^2$, which is generically not $1/N$.

The gap is structural:

\begin{quote}
Admissibility is a binary filter. Probability requires a measure. The filter determines the support; the measure determines the weight.
\end{quote}

What is needed is a principled construction of a state-dependent weight $w_i(\Psi, \Ocal)$ for each admissible sector, such that:
\begin{equation}\label{eq:weight}
    P(\Phi_i \mid \Psi, \Ocal) = \frac{w_i(\Psi, \Ocal)}{\sum_{j=1}^N w_j(\Psi, \Ocal)}
\end{equation}

The remainder of this paper develops five routes to $w_i$. Routes~A and~B are \emph{motivational prototypes}: they illustrate why probability must be geometric in character, but neither produces quadratic weighting by itself, and neither is advanced here as a complete candidate. Route~C is the paper's main dynamical derivation: under the equilibrium, sector-supported, sampling, and state-completeness hypotheses of Proposition~\ref{prop:born}, the stationary measure of the closure dynamics yields sector masses reproducing the Born form $|c_i|^2$. Route~D is a Gleason-style structural confirmation once Hilbert-space structure is assumed. Route~E assembles prior-paper inputs (the Paper~XVIII complex-structure bookkeeping and the Paper~XXI equilibrium-tangent $U(N)$-invariant K\"ahler structure on $\mathbb{CP}^{N-1}$ under its sector-subspace invariance hypothesis) together with a local Laplace identification of the closure Hessian with a configuration-space Fisher metric, and invokes Gleason's frame-function theorem to identify the Born measure as the unique assignment satisfying (U1)--(U4) and universal completeness. The Chentsov--Petz content provides a structural cross-check (the imported $U(N)$-invariant K\"ahler structure is Fubini--Study up to scale) rather than an independent uniqueness driver. Routes~C (existence) and~E (uniqueness) carry the central argument; Routes~A, B, and~D are preparatory or confirmatory.

%==================================================================
\section{Route A: Basin-Volume Weighting}\label{sec:basin}
%==================================================================

\begin{definition}[Basin of attraction]\label{def:basin}
For a locally stable configuration $\Phi_i^*$ (local minimum of $\Fcal$ in $\Phi_{\mathrm{adm}}$), the \textit{basin of attraction} is:
\begin{equation}
    B_i = \bigl\{\, \Phi \in \mathbb{R}^\Ecal \;\big|\; \text{gradient flow of } \Fcal \text{ starting at } \Phi \text{ converges to } \Phi_i^* \,\bigr\}
\end{equation}
\end{definition}

\begin{definition}[Basin-volume weight]\label{def:basin_weight}
\begin{equation}\label{eq:basin_weight}
    w_i^{(\mathrm{A})} = \mu(B_i \cap \Phi_{\mathrm{adm}})
\end{equation}
where $\mu$ is a natural measure on configuration space (e.g.\ Lebesgue measure on the finite-dimensional event set).
\end{definition}

\textbf{Interpretation.} If the system's initial configuration is drawn from a distribution that is approximately uniform over the admissible region, then the probability of ending in sector $i$ is proportional to the volume of its basin. Larger basins attract more initial conditions.

\begin{remark}[Strengths and limitations]
Basin-volume weighting is natural for classical dynamical systems. It explains why different outcomes have different probabilities without invoking randomness: the asymmetry comes from geometry. However, it does not by itself explain why the weights should be \textit{quadratic} in any amplitude. It produces a probability measure, but not necessarily the Born measure.
\end{remark}

%==================================================================
\section{Route B: Stability-Depth Weighting}\label{sec:stability}
%==================================================================

\begin{definition}[Stability-depth weight]\label{def:stability_weight}
\begin{equation}\label{eq:stability_weight}
    w_i^{(\mathrm{B})} = e^{-\beta\, \Fcal[\Phi_i^*]}
\end{equation}
where $\Fcal[\Phi_i^*]$ is the closure-functional value at the $i$-th minimum and $\beta > 0$ is a scale parameter.
\end{definition}

\textbf{Interpretation.} Configurations with lower closure cost (deeper minima) are more heavily weighted. This is the Boltzmann-like weighting familiar from statistical mechanics: outcomes that are more constraint-compatible are more probable.

\begin{remark}
This route connects to the thermodynamic structure of the closure landscape. It produces non-equal probabilities whenever the minima have different depths. But the Boltzmann form $e^{-\beta E}$ is not the same as the Born form $|\psi|^2$ unless $\beta$ and $\Fcal$ are related to amplitude in a specific way. Route~C makes this connection explicit.
\end{remark}

%==================================================================
\section{Route C: Stationary-Measure Inheritance}\label{sec:stationary}
%==================================================================

This is the strongest route, because it connects directly to established results from the programme.

\subsection{The Paper~XIV stationary distribution}

Paper~XIV~\citep{SmartXIV} established that the stochastic closure dynamics (Langevin equation $d\psi = -\nabla\Fcal\, dt + \sigma\, dW$) admits a unique stationary distribution:
\begin{equation}\label{eq:stationary}
    \Pst(\psi) \propto e^{-2\Fcal(\psi)/\sigma^2}
\end{equation}
This configuration-space distribution concentrates on the low-$\Fcal$ regions of the constraint manifold in configuration space (the closure attractor basins of Paper~XIV; not to be confused with the projective sector manifold $\Mcl(\Ocal) \subseteq \mathbb{CP}^{N-1}$ introduced in Definition~\ref{def:closure_manifold}), and provides the probabilistic structure of the closure landscape.

\subsection{The Nelson wavefunction connection}

Paper~XVIII~\citep{SmartXVIII} constructed the Nelson wavefunction $\Psi = \sqrt{\rho}\, e^{iS/\sigma^2}$ from the paired forward/backward closure processes. At the Paper~XIV equilibrium, the density $\rho$ equals $\Pst$, so:
\begin{equation}\label{eq:born_connection}
    |\Psi|^2 = \rho = \Pst \propto e^{-2\Fcal/\sigma^2}
\end{equation}

This is the central observation: \textit{in the equilibrium regime, the effective Born density $|\Psi|^2$ coincides with the stationary measure of the closure dynamics.}

\subsection{Why measurement samples the stationary sector measure}

The connection between the stationary measure and measurement outcomes rests on the following reasoning:
\begin{enumerate}
    \item The underlying stochastic closure dynamics has a unique stationary measure $\Pst$ (Paper~XIV).
    \item Observer constraints carve this measure into admissible sectors (Paper~XXIX).
    \item Measurement frequencies correspond to repeated sampling from the constrained stationary ensemble: each measurement draws from the same $\Pst$, restricted to the admissible set.
\end{enumerate}
This paper treats this reasoning as a working hypothesis rather than a completed theorem. In particular, the assumption that the system has equilibrated to $\Pst$ at the time of measurement is not proved for general measurement contexts.

\subsection{Stationary-measure weight}

\begin{definition}[Stationary-measure weight]\label{def:stationary_weight}
\begin{equation}\label{eq:stationary_weight}
    w_i^{(\mathrm{C})} = \int_{\Phi_i} \Pst(\Phi)\, d\mu(\Phi) = \int_{\Phi_i} |\Psi(\Phi)|^2\, d\mu(\Phi)
\end{equation}
the total stationary probability mass in outcome sector $\Phi_i$.
\end{definition}

\begin{proposition}[Conditional Born weighting from stationary sector measure]\label{prop:born}
Assume all of the following:
\begin{enumerate}
    \item The system is at the Paper~XIV equilibrium ($\rho = \Pst$). The proof establishes the formula below only at equilibrium; a near-equilibrium version would carry an additional $O(\rho - \Pst)$ error term not tracked here.
    \item The outcome sectors $\Phi_i$ are well-separated basins (negligible overlap).
    \item The Nelson wavefunction $\Psi$ admits a decomposition $\Psi = \sum_i c_i \phi_i$ where each $\phi_i$ is normalised ($\|\phi_i\| = 1$) and supported on $\Phi_i$.
    \item \textbf{(Sampling hypothesis.)} Measurement outcome frequencies equal the integrated stationary mass over sectors, normalised over the observer-admissible set: $P(\Phi_i \mid \Ocal) = \int_{\Phi_i}|\Psi|^2\,d\mu / \int_{\Phi_\mathrm{adm}}|\Psi|^2\,d\mu$. This is an additional measurement-theoretic postulate, not a consequence of the stationary measure itself.
    \item \textbf{(State-completeness.)} The Nelson density $|\Psi|^2$ is supported within $\Phi_\mathrm{adm}(\Ocal)$, so $\int_{\Phi_\mathrm{adm}}|\Psi|^2 \, d\mu = \|\Psi\|^2 = 1$. This is the state-completeness hypothesis of Remark~\ref{rem:complete_partial}; without it, the $|c_i|^2$ identification below holds only up to the observer-dependent normalisation $\sum_{j \in \mathrm{adm}}|c_j|^2$.
\end{enumerate}
Then, under all five hypotheses:
\begin{equation}\label{eq:born_result}
    P(\Phi_i \mid \Ocal) \;=\; \frac{\int_{\Phi_i} |\Psi|^2\, d\mu}{\int_{\Phi_{\mathrm{adm}}} |\Psi|^2\, d\mu} \;=\; |c_i|^2 + O(\text{overlap}),
\end{equation}
which reduces to the standard Born-rule form
\begin{equation}\label{eq:born_exact}
    P(\Phi_i \mid \Ocal) \;\approx\; |c_i|^2 \qquad \text{(equal in the disjoint-support limit).}
\end{equation}
Without hypothesis (5) (partial observer), the sampling ratio instead gives the conditional form $|c_i|^2 / \sum_{j \in \mathrm{adm}} |c_j|^2$ --- see Remarks~\ref{rem:partial}, \ref{rem:vs_standard_qm}.
\end{proposition}

\begin{proof}
This is a sector-mass decomposition lemma plus the sampling and completeness hypotheses. Under (3), $|\Psi|^2 = |\sum_i c_i \phi_i|^2$. Under (2) and the normalised-$\phi_i$ clause of (3), the supports of $\phi_i$ are effectively disjoint, so cross terms $c_i^* c_j \phi_i^* \phi_j$ are negligible when integrated over any single sector $\Phi_i$; the correction scales with the overlap and vanishes in the disjoint-support limit. Therefore $\int_{\Phi_i}|\Psi|^2 = |c_i|^2 + O(\text{overlap})$. Under (4), the measurement frequencies equal the normalised ratio of sector masses; under (5), the denominator equals $\|\Psi\|^2 = 1$, so the ratio reduces to $|c_i|^2 + O(\text{overlap})$. Under (1), the Nelson density is the stationary density, so the integrated masses are integrals of $\Pst$.
\end{proof}

\begin{remark}[The role of the sampling hypothesis (4)]\label{rem:sampling}
Assumption (4) is the step that connects sector-mass integrals to measurement frequencies. It is \emph{not} inherited from the stationary measure of the closure dynamics: the stationary measure is a property of the dynamics, not of the measurement procedure. A complete derivation would need to show that measurement interactions in the closure framework do sample the constrained stationary ensemble --- which would require formalising ergodicity and mixing properties of the closure Markov process, together with a treatment of decoherence. Proposition~\ref{prop:born} is therefore a sector-mass decomposition lemma \emph{plus} the sampling postulate; it is not, by itself, a derivation of measurement statistics from closure alone.
\end{remark}

\begin{remark}[What this does and does not establish]
Proposition~\ref{prop:born} shows that \textit{if} the system is at equilibrium ($\rho = \Pst$; see Section~\ref{sec:discussion}'s subsection on the equilibrium assumption), \textit{if} outcome sectors are well-separated, \textit{if} the wavefunction decomposes into sector-supported components, \textit{if} measurement samples the constrained stationary ensemble (assumption~4), \emph{and} \textit{if} the Nelson density is supported in the observer-admissible set (state-completeness, assumption~5), then the Born rule follows from the stationary measure of the closure dynamics. The first assumption (equilibrium) concerns the dynamical regime of Papers~XVIII--XXI. Assumptions (2) and (3) --- well-separated sectors and a sector-supported decomposition --- are additional modelling assumptions of the present paper; they are \emph{not} proved in Paper~XVIII or Paper~XXI, which provide the Nelson pairing and the equilibrium-tangent projective structure but do not by themselves construct sector-localised basis vectors or well-separated sector boundaries. Assumption (4) is a measurement-theoretic postulate; assumption (5) is the observer-side completeness condition, without which the ratio reduces to the observer-conditional form $|c_i|^2 / \sum_{j\in\mathrm{adm}}|c_j|^2$ rather than $|c_i|^2$. Whether all five are satisfied in physical measurement contexts is an open question.
\end{remark}

\begin{remark}[Logical status]
Within the equilibrium closure regime, and after the Route~C sampling postulate and state-completeness assumption are imposed, the quadratic probability law need not be introduced as an independent measurement axiom; it is inherited from the stationary measure. The chain is:
\begin{enumerate}
    \item the stationary distribution $\Pst \propto e^{-2\Fcal/\sigma^2}$ (derived in Paper~XIV),
    \item the identification $|\Psi|^2 = \Pst$ at equilibrium (derived in Paper~XVIII),
    \item the decomposition of $\Pst$ into outcome-sector integrals (Definition~\ref{def:stationary_weight}),
    \item the sampling postulate (Proposition~\ref{prop:born} hypothesis (4)) and state-completeness (hypothesis (5)), which together bridge stationary sector masses to measurement frequencies and the unconditional $|c_i|^2$ form.
\end{enumerate}
The quadratic form is therefore not imposed independently, but it is not extracted from the stationary measure alone either: the sampling and completeness assumptions are indispensable. Whether this inheritance extends to all measurement contexts, including those far from equilibrium, is an open question.
\end{remark}

\begin{remark}[Three levels of identification]
Three levels should be distinguished. First, $\Pst$ is the exact stationary measure of the stochastic closure process (Paper~XIV). Second, in the equilibrium tangent regime of Papers~XVIII--XXI, this stationary density is represented effectively as $|\Psi|^2$. Third, measurement frequencies are hypothesised to sample the induced sector masses of that equilibrium measure. The present paper argues that, under these identifications and assumptions, the observed quadratic law is inherited rather than postulated.
\end{remark}

%==================================================================
\section{Route D: Gleason-Type Argument}\label{sec:gleason}
%==================================================================

If the admissible-sector geometry has Hilbert-space structure, an independent argument constrains the form of the probability measure.

\subsection{The Gleason bridge}

Gleason's theorem~\citep{Gleason1957} states: on a Hilbert space $\mathcal{H}$ of dimension $\geq 3$, the only non-contextual, countably additive probability measure $\mu$ on the lattice of closed subspaces satisfies $\mu(P) = \mathrm{tr}(\rho P)$ for some density operator $\rho$. For a pure state $\rho = |\psi\rangle\langle\psi|$, this gives:
\begin{equation}
    P_i = \mathrm{tr}(|\psi\rangle\langle\psi| \cdot \Pi_i) = |\langle \phi_i | \psi \rangle|^2
\end{equation}
the Born rule.

\subsection{Application to the VFD framework}

Papers~XVIII--XXI established that the equilibrium tangent dynamics of the closure framework is the Schr\"odinger equation, which naturally lives on a Hilbert space. If the observer-admissible outcome sectors correspond to orthogonal subspaces of this effective Hilbert space, and if the probability measure over sectors is required to be non-contextual and additive, then Gleason's theorem implies the Born rule.

\begin{remark}[What VFD adds to Gleason]
Gleason's theorem assumes Hilbert-space structure and derives the Born rule. Route~D does not provide the physics; it provides a \textit{uniqueness confirmation}. The VFD contribution is to explain \textit{why} the relevant space has Hilbert structure: because the equilibrium tangent dynamics is Schr\"odinger-type (Papers~XVIII--XXI). The Gleason argument then confirms that, given this structure, the Born rule is the only compatible measure. The physical origin of the weighting comes from Route~C (the stationary measure); Route~D establishes that no other weighting is consistent with the emergent Hilbert structure.
\end{remark}

\begin{remark}[Assumptions]
This route requires: (1) the system is in the equilibrium tangent regime; (2) the outcome sectors correspond to orthogonal subspaces; (3) the dimension is $\geq 3$; (4) the probability measure is non-contextual and additive. These are the standard Gleason assumptions. Whether the VFD framework satisfies them in full generality is an open question.
\end{remark}

%==================================================================
\section{Route E: Imported Geometric Packaging for a Gleason Uniqueness Argument}\label{sec:chentsov}
%==================================================================

Routes~A--D establish that quadratic weighting is \emph{compatible} with the closure framework. This section assembles the geometric inputs on which a Gleason-style uniqueness argument becomes applicable: a closure-side local identification of the Hessian with a configuration-space Fisher metric, and a $U(N)$-invariant K\"ahler structure on $\mathbb{CP}^{N-1}$ imported from Paper~XXI as standard geometric QM on the assumed invariant subspace. Under the conditions stated in (U1)--(U4) below (Paper~XXI supplies the equilibrium-tangent unitary/projective geometric action, with covariance of $P$ assumed as hypothesis (U1), non-contextuality taken as a hypothesis with a restricted form available from Paper~XXIX, continuity, and eigenstate certainty), together with universal completeness $\Mcl(\Ocal) = \mathbb{CP}^{N-1}$ and basis-realisability (Remark~\ref{rem:basis_realisability}), the Born measure is the unique probability assignment satisfying these hypotheses for $N \geq 3$. The $N = 2$ case is not covered here (Remark~\ref{rem:N_two}). The uniqueness proof itself is Gleason's; the role of the information-geometric content is structural (Petz identifies the imported $U(N)$-invariant K\"ahler structure as Fubini--Study up to scale) rather than load-bearing for the Born derivation.

The argument combines three elements: (i) a \emph{local, Laplace-regime} identification of the closure Hessian with the Fisher metric of a location-parameter family on configuration space (Proposition~\ref{prop:hessian_fisher}) --- this is closure-side; (ii) a \emph{global, $U(N)$-invariant} K\"ahler structure on $\mathbb{CP}^{N-1}$ --- this is imported from Paper~XXI's equilibrium-tangent construction; (iii) Gleason's frame-function theorem~\citep{Gleason1957,Peres1995,CookKeaneMoran1985} applied to the combined structure. Petz's theorem is used only as a structural cross-check identifying the imported $U(N)$-invariant K\"ahler structure as Fubini--Study up to scale; it does not enter the uniqueness proof.

Where Route~D imports Hilbert structure as an unexplained assumption, Route~E combines a local closure-side metric identification with the global equilibrium-tangent K\"ahler structure of Paper~XXI. The result is not a derivation of the global K\"ahler structure from closure alone; it is a chain that says: \emph{given} Papers~XVIII, XXI, XXIX, the Laplace regime, universal completeness $\Mcl(\Ocal) = \mathbb{CP}^{N-1}$, basis-realisability (Remark~\ref{rem:basis_realisability}), and axioms (U1)--(U4) on a global probability assignment $P$ on $\mathbb{CP}^{N-1}$, Born weighting is the unique assignment so constrained for $N \geq 3$. Section~\ref{subsec:route_e_delta} makes the external inputs fully explicit.

The logical structure, factored into its dependency chain: (i) Paper~XXI supplies a $U(N)$-invariant K\"ahler structure on $\mathbb{CP}^{N-1}$ at the equilibrium tangent, as imported standard geometric QM on the assumed invariant subspace; (ii) Paper~XVIII supplies the complex structure as a bookkeeping observation on the Nelson wavefunction, and Proposition~\ref{prop:hessian_fisher} supplies at most a candidate local Riemannian normalisation via the closure Hessian in the Laplace regime, conditional on an unproved bridge theorem (Remark~\ref{rem:global_gap}); (iii) Petz's theorem identifies the imported $U(N)$-invariant K\"ahler structure with Fubini--Study up to scale; (iv) Gleason's frame-function theorem, applied to $U(N)$-covariant, non-contextual, continuous, eigenstate-certain probability assignments in dimension $N \geq 3$, fixes the unique compatible probability as the Born measure under universal completeness and basis-realisability. The following subsections execute this chain step by step.

\subsection{The admissible set as a statistical parameter space}\label{subsec:stat_manifold}

Let $\Phi_{\mathrm{adm}}(\Ocal)$ denote the observer-admissible configuration set (Definition~\ref{def:admissible}). In the equilibrium tangent regime of Papers~XVIII--XXI, each admissible outcome sector $\Phi_i$ is associated with a normalised sector wavefunction $\phi_i$, and a general admissible state is the superposition
\begin{equation}
    \Psi = \sum_{i=1}^{N} c_i\, \phi_i, \qquad \sum_{i=1}^{N} |c_i|^2 = 1,
\end{equation}
parameterised by complex coefficients $(c_1, \ldots, c_N)$ modulo global phase. This gives the finite-dimensional complex projective space $\mathbb{CP}^{N-1}$.

\begin{definition}[Closure statistical parameter space]\label{def:closure_manifold}
Recall that each $\Psi \in \Hcal_\Ocal$ induces, via the Nelson pairing of Paper~XVIII, a density $\rho_\Psi = |\Psi|^2$ on configuration space $\mathbb{R}^\Ecal$. Define the admissible projective subset
\begin{equation}
    \Mcl(\Ocal) \;=\; \bigl\{\, [\Psi] \in \mathbb{CP}^{N-1} \;\big|\; \mathrm{supp}(\rho_\Psi) \subseteq \Phi_{\mathrm{adm}}(\Ocal) \,\bigr\},
\end{equation}
i.e.\ the set of projective states whose Nelson-induced density is concentrated on the observer-admissible configurations. $\Mcl(\Ocal)$ is equipped with the family $\{P_i(\Psi)\}_{i=1}^{N}$ of sector-probability assignments obtained by integrating $\rho_\Psi$ over each sector (equation~\eqref{eq:prob_via_integration}).
\end{definition}

\begin{remark}[Typing]\label{rem:typing}
The definition typechecks: $[\Psi] \in \mathbb{CP}^{N-1}$ is a Hilbert-space object; $\Phi_{\mathrm{adm}}(\Ocal) \subset \mathbb{R}^\Ecal$ is a configuration-space object. The two are connected by the Nelson pairing $\Psi \mapsto \rho_\Psi$, which is a map $\Hcal_\Ocal \to L^1(\mathbb{R}^\Ecal)$. Membership in $\Mcl(\Ocal)$ is formulated as a support condition on $\rho_\Psi$, not as direct membership $\Psi \in \Phi_{\mathrm{adm}}$ (which would be ill-typed). In the sector-supported decomposition regime assumed by Route~C, $\mathrm{supp}(\rho_\Psi) \subseteq \bigsqcup_i \Phi_i(\Ocal) \subseteq \Phi_{\mathrm{adm}}(\Ocal)$ automatically.
\end{remark}

Each point of $\Mcl(\Ocal)$ parameterises a probability distribution over outcome sectors via equation~\eqref{eq:prob_via_integration}. $\Mcl(\Ocal)$ can therefore be treated formally as a statistical parameter space in the sense of information geometry~\citep{AmariNagaoka2000}: the parameter space of a formally smooth family of probability distributions over a common outcome set. Proving genuine smooth-manifold structure and smoothness of the probability family is beyond the scope of the present paper.

\subsection{Covariance groups for the probability assignment}\label{subsec:sufficient}

The forcing argument below requires that the sector-probability assignment $P(\Pi, \Psi)$ be covariant under a specified group of unitary transformations on $\Hcal_\Ocal$. Two groups appear, and they are \emph{not} the same:

\begin{enumerate}
    \item \textbf{Sector-preserving unitaries} (the narrow group). These permute the basis vectors $\phi_i \in \Hcal_\Ocal$ up to phases. Formally, $T \in U(\Hcal_\Ocal)$ is sector-preserving if there exist a permutation $f$ of $\{1, \ldots, N\}$ and phases $e^{i\alpha_i}$ with
    \begin{equation}\label{eq:cp_unitary}
        T \phi_i \;=\; e^{i\alpha_i}\, \phi_{f(i)} \qquad \text{for every sector basis vector } \phi_i.
    \end{equation}
    This group is strictly smaller than $U(N)$: it is the semidirect product $U(1)^N \rtimes S_N$ of diagonal phases and sector permutations.
    \item \textbf{Equilibrium-tangent unitaries} (the full group). Paper~XXI~\citep{SmartXXI} establishes that the linearised Schr\"odinger dynamics at the equilibrium tangent is unitary on $\Hcal_\Ocal$ and $U(N)$-covariant: every $U \in U(N)$ acts as a K\"ahler isometry on the equilibrium-tangent structure of $\mathbb{CP}^{N-1}$. A generic element of $U(N)$ sends basis vectors to superpositions and therefore does \emph{not} satisfy~\eqref{eq:cp_unitary}; it is not sector-preserving.
\end{enumerate}

\begin{definition}[Sector-preserving unitary]\label{def:cp_unitary}
A unitary $T \in U(\Hcal_\Ocal)$ is \emph{sector-preserving} if it satisfies~\eqref{eq:cp_unitary} for some index permutation $f$ and phases $\{\alpha_i\}$.
\end{definition}

\begin{proposition}[Sector-preserving unitaries induce sector-probability covariance, under the full Route~C hypotheses]\label{prop:suff}
Adopt all five hypotheses of Proposition~\ref{prop:born} (equilibrium, well-separated sectors, sector-supported decomposition, the sampling hypothesis, and state-completeness) for both $\Psi$ and $T\Psi$. Let $T \in U(\Hcal_\Ocal)$ be a sector-preserving unitary with index permutation $f$ (Definition~\ref{def:cp_unitary}). Let $\Psi \in \Mcl(\Ocal)$ admit a sector-supported decomposition $\Psi = \sum_j c_j \phi_j$ with $|\phi_j|^2$ supported on $\Phi_j(\Ocal)$; state-completeness for $\Psi$ transfers automatically to $T\Psi$ if the sectors cover the admissible set and $\sum_i|c_i|^2 = 1$. Then the Route~C sector-probability family $\{P_i\}$ satisfies
\begin{equation}\label{eq:sufficiency}
    P_{f(i)}(T\Psi) = P_i(\Psi) \qquad \forall i \in \{1, \ldots, N\}.
\end{equation}
Without the state-completeness assumption, covariance holds instead for the unnormalised sector mass $|c_i|^2$ (numerator of equation~\eqref{eq:prob_via_integration}); the conditional probability can fail to be covariant if the admissible denominator differs between $\Psi$ and $T\Psi$.
\end{proposition}

\begin{proof}
By Proposition~\ref{prop:born} (all five hypotheses, including state-completeness), the sector-probability assignment on a sector-supported state reduces to $P_i(\Psi) = |c_i|^2$. Write $\Psi = \sum_j c_j \phi_j$. By~\eqref{eq:cp_unitary},
\begin{equation}
    T\Psi \;=\; \sum_j c_j T\phi_j \;=\; \sum_j c_j\, e^{i\alpha_j}\, \phi_{f(j)} \;=\; \sum_k \bigl(c_{f^{-1}(k)}\, e^{i\alpha_{f^{-1}(k)}}\bigr) \phi_k,
\end{equation}
so the coefficients of $T\Psi$ in the sector basis are $c'_k = c_{f^{-1}(k)} e^{i\alpha_{f^{-1}(k)}}$. Therefore $|c'_k|^2 = |c_{f^{-1}(k)}|^2$, and since permutations preserve total amplitude $\sum_k |c'_k|^2 = \sum_j |c_j|^2$, state-completeness for $\Psi$ transfers to $T\Psi$. Applying Proposition~\ref{prop:born} again to $T\Psi$ gives $P_{f(i)}(T\Psi) = |c'_{f(i)}|^2 = |c_{f^{-1}(f(i))}|^2 = |c_i|^2 = P_i(\Psi)$, which is~\eqref{eq:sufficiency}.

For the partial-observer case (without state-completeness), the denominator of equation~\eqref{eq:prob_via_integration} is $\sum_{j\in\mathrm{adm}}|c_j|^2$ for $\Psi$ and $\sum_{k\in\mathrm{adm}}|c'_k|^2 = \sum_{j\in f^{-1}(\mathrm{adm})}|c_j|^2$ for $T\Psi$; these need not agree unless the admissibility pattern is $f$-invariant, and so the conditional ratio is not covariant in general.
\end{proof}

\begin{remark}[Why this proposition is not enough by itself]\label{rem:suff_scope}
Proposition~\ref{prop:suff} only delivers covariance under the narrow group of sector-preserving unitaries. Gleason's argument in Theorem~\ref{thm:born_unique}, Step~1, relies on covariance under the \emph{full} $U(N)$ action: the reduction of $P(\Pi, \Psi)$ to a function of the scalar invariant $\langle\Psi|\Pi|\Psi\rangle$ uses that every pair $(\Pi, \Psi)$ with fixed inner product lies on a single $U(N)$-orbit, which requires $U(N)$-equivariance of $P$, not merely sector-preservation. The full $U(N)$-covariance is \emph{not} supplied by Proposition~\ref{prop:suff}; it is supplied by Paper~XXI's equilibrium-tangent construction (see the next paragraph and axiom (U1) of Theorem~\ref{thm:born_unique}).
\end{remark}

\textbf{Source of full $U(N)$-covariance.} Paper~XXI~\citep{SmartXXI} establishes that, at the equilibrium tangent, the linearised Schr\"odinger dynamics is generated by a self-adjoint Hamiltonian and preserves the Hermitian inner product on $\Hcal_\Ocal$; consequently the full unitary group $U(N)$ acts as K\"ahler isometries of the equilibrium-tangent structure on $\mathbb{CP}^{N-1}$. This $U(N)$-invariance is the equilibrium-tangent fact imported as axiom (U1) of Theorem~\ref{thm:born_unique}. The present paper does not derive $U(N)$-invariance from the full nonlinear closure dynamics; the invariance is a property of the linearised regime identified in Paper~XXI.

The \emph{sector-preserving} subgroup (Proposition~\ref{prop:suff}) and the \emph{full} $U(N)$ action (Paper~XXI) play different roles: the former gives finite index-permutation symmetry of the probability assignment; the latter gives continuous equivariance needed for Gleason's frame-function reduction. Both are used below, and we keep them notationally distinct.

\subsection{Chentsov's uniqueness theorem}\label{subsec:chentsov_theorem}

\begin{theorem}[Chentsov 1972~\citep{Chentsov1982}; Petz quantum extension~\citep{Petz1996} (imported; not proved here)]\label{thm:chentsov}
\begin{enumerate}
    \item (\emph{Classical.}) On a finite-dimensional statistical manifold of dimension $\geq 2$, the Fisher--Rao metric is the unique Riemannian metric (up to constant positive multiple) invariant under sufficient statistics (equivalently, Markov morphisms).
    \item (\emph{Standard geometric-QM fact, used here for structural comparison only.}) On the complex projective manifold $\mathbb{CP}^{N-1}$ with $N \geq 2$, the $U(N)$-invariant Fubini--Study K\"ahler structure is determined up to a positive scalar. In the present paper this is used only to identify the imported $U(N)$-invariant K\"ahler structure of Paper~XXI with Fubini--Study up to scale, as a structural cross-check; no new quantum uniqueness result is proved here.
    \item (\emph{Real-slice bridge.}) Let $\mathbb{RP}^{N-1}_+ \subset \mathbb{CP}^{N-1}$ denote the real, non-negative-coefficient slice, parameterised by $(c_1, \ldots, c_N) \in \mathbb{R}^N_{\geq 0}$ with $\sum_i c_i^2 = 1$. The modulus-squared sector projection
    \begin{equation}
        \mu: \mathbb{RP}^{N-1}_+ \to \Delta^{N-1}, \qquad (c_1, \ldots, c_N) \mapsto (c_1^2, \ldots, c_N^2)
    \end{equation}
    pushes the \emph{restriction} of the Fubini--Study metric $g_{\mathrm{FS}}$ to $\mathbb{RP}^{N-1}_+$ forward to a constant multiple of the Fisher--Rao metric $g_{\mathrm{FR}} = \sum_i (dp_i)^2/p_i$ on the interior of the simplex $\Delta^{N-1}$. Equivalently, under $c_i = \sqrt{p_i}$, the real-slice metric $\sum_i dc_i^2$ pulls back to $\tfrac14 g_{\mathrm{FR}}$ on $\Delta^{N-1}$.
\end{enumerate}
\end{theorem}

Parts (1) and (2) are standard (see~\citep{AmariNagaoka2000,Petz1996}). Part (3) is the classical identification of the real-slice Fubini--Study metric with (a scaling of) Fisher--Rao; concretely, $(dc_i)^2 = (dp_i)^2/(4 p_i)$, so $g_{\mathrm{FS}}|_{\mathbb{RP}^{N-1}_+} = \tfrac14 g_{\mathrm{FR}}$ after modulus-squared projection. Note: Part (3) is a real-slice statement. A general pure state in $\mathbb{CP}^{N-1}$ has complex coefficients $c_i$ whose relative phases are lost under $\mu(\Psi) = (|c_i|^2)$, so $\mu$ is not an isometric embedding of the full $\mathbb{CP}^{N-1}$ into $\Delta^{N-1}$: the push-forward collapses the phase degrees of freedom. Within Route~E, Part~(3) is used only to identify the Riemannian scale of the real-slice tangent at the equilibrium state (Proposition~\ref{prop:hessian_fisher}), not as a full pushforward between the complex projective manifold and the probability simplex.

\subsection{The closure Hessian induces the local Fisher metric}\label{subsec:hessian_fisher}

The Fisher structure on $\Mcl$ is not imposed from outside. It arises from the second-order structure of the closure functional at equilibrium --- but as a \emph{local} identification in the Laplace (Gaussian) regime around an equilibrium point, not a global identity.

\begin{proposition}[Closure Hessian as the Fisher metric of a location-parameter family on configuration space, in the Laplace regime]\label{prop:hessian_fisher}
Let $\phi^* \in \mathbb{R}^\Ecal$ be a non-degenerate local minimum of $\Fcal$ (a configuration-space equilibrium point, \emph{not} a projective state). Consider the \emph{location-parameter family} of location-translated densities on configuration space (each element is the stationary distribution of a closure functional translated by $\theta$, not the stationary distribution of the original $\Fcal$), indexed by a translation parameter $\theta \in \mathbb{R}^d$ with $d = \dim(\mathbb{R}^\Ecal) = |\Ecal|$:
\begin{equation}
    \Pst(\Phi;\theta) \;=\; \frac{1}{Z(\theta)}\, e^{-2\Fcal(\Phi - \phi^*(\theta))/\sigma^2},
    \qquad \phi^*(\theta) \;=\; \phi^* + \theta,
    \qquad Z(\theta) = \int_{\mathbb{R}^\Ecal} e^{-2\Fcal(\Phi - \phi^*(\theta))/\sigma^2}\,d\mu(\Phi).
\end{equation}
Each element of the family is a probability distribution on the configuration space $\mathbb{R}^\Ecal$; the family is parameterised by translations $\theta$ of the equilibrium point in $\mathbb{R}^\Ecal$. In the Laplace (small-$\sigma$) approximation where $\Pst$ is well-approximated by a Gaussian of inverse covariance $(2/\sigma^2) H[\phi^*]$, the Fisher information matrix of this family at $\theta = 0$ is
\begin{equation}\label{eq:hessian_fisher}
    g^{\mathrm{Fisher}}_{ab}(0)
    \;=\; \mathbb{E}_{\Pst}\!\left[\partial_a \log\Pst \cdot \partial_b \log\Pst\right]
    \;=\; \frac{2}{\sigma^2}\, H_{ab}[\phi^*] \;+\; O(\sigma^0),
\end{equation}
where $H_{ab}[\phi^*] = \partial_a\partial_b \Fcal|_{\phi^*}$ is the closure Hessian. This is a metric on the location-parameter family (a $d$-dimensional submanifold of the space of probability densities on $\mathbb{R}^\Ecal$); it is \emph{not} a metric on $\mathbb{CP}^{N-1}$.
\end{proposition}

\begin{proof}
Differentiating $\log\Pst = -2\Fcal(\Phi - \phi^*(\theta))/\sigma^2 - \log Z(\theta)$ with respect to $\theta^a$ gives the score
\begin{equation}
    \partial_a \log\Pst(\Phi;\theta)
    = \frac{2}{\sigma^2}\,\partial_a \Fcal(\Phi - \phi^*(\theta))
    \;-\; \partial_a \log Z(\theta).
\end{equation}
The Fisher information is the covariance of the score; its expectation vanishes, and the $\log Z$ term is deterministic (it drops out of the covariance):
\begin{equation}
    g^{\mathrm{Fisher}}_{ab}(\theta)
    \;=\; \frac{4}{\sigma^4}\, \mathrm{Cov}_{\Pst}\!\bigl[\partial_a \Fcal,\, \partial_b \Fcal\bigr].
\end{equation}
In the Laplace regime, $\Pst$ is approximately Gaussian with covariance $C = (\sigma^2/2)\,H^{-1}[\phi^*]$. Writing $\partial_a\Fcal(\Phi - \phi^*) \approx H_{ac}(\Phi - \phi^*)^c$ to leading order, we obtain $\mathrm{Cov}[\partial_a\Fcal, \partial_b\Fcal] \approx H_{ac} C^{cd} H_{db} = (\sigma^2/2)\,H_{ab}$. Substituting,
\begin{equation}
    g^{\mathrm{Fisher}}_{ab}(0) \;\approx\; \frac{4}{\sigma^4} \cdot \frac{\sigma^2}{2}\, H_{ab}[\phi^*] \;=\; \frac{2}{\sigma^2}\, H_{ab}[\phi^*],
\end{equation}
with correction terms of order $\sigma^0$ controlled by higher derivatives of $\Fcal$.
\end{proof}

\begin{remark}[Scope of the identification and how it connects to $\mathbb{CP}^{N-1}$]\label{rem:hessian_scope}
Proposition~\ref{prop:hessian_fisher} identifies the Fisher information of the location-parameter family on configuration space $\mathbb{R}^\Ecal$ with the closure Hessian; it is a result about a $d$-dimensional submanifold of probability densities, not about $\mathbb{CP}^{N-1}$. A conjectural connection to the projective state manifold would go via Paper~XVIII's Nelson pairing: \emph{if} each location-translated stationary density can be realised as an admissible state $\Psi \in \Hcal_\Ocal$ with $|\Psi|^2 = \Pst(\cdot;\theta)$ (a bridge assumption not proved in Paper~XVIII, XXI, or the present paper), then the tangent direction $\partial_\theta \Psi$ lies in the real slice of $\Hcal_\Ocal$. Under that bridge assumption, the Fisher metric on the location-parameter family pulls back to a metric on a $d$-dimensional submanifold of $\mathbb{CP}^{N-1}$ (the real slice containing $[\Psi^*]$). The full K\"ahler metric on $\mathbb{CP}^{N-1}$ is \emph{not} produced by this proposition alone; it is supplied by Paper~XXI's equilibrium-tangent $U(N)$-invariant construction, and Proposition~\ref{prop:fs_assembly} uses~\eqref{eq:hessian_fisher} only to supply a \emph{candidate} overall normalisation at the equilibrium point, conditional on the bridge being established. Away from equilibrium, or at higher order in $\sigma$, the Laplace approximation breaks down and the Fisher--Hessian identification is no longer exact.
\end{remark}

\subsection{The K\"ahler structure on the admissible parameter space}\label{subsec:kahler}

The present paper does not establish a K\"ahler structure on $\Mcl(\Ocal)$ from closure alone. It imports the standard $U(N)$-invariant K\"ahler structure on $\mathbb{CP}^{N-1}$ from Paper~XXI (under its sector-subspace invariance hypothesis) and compares it locally to closure-side Fisher data. The three pieces of that imported K\"ahler structure enter via separate prior papers:

\begin{enumerate}
    \item \textbf{Riemannian.} The local Fisher metric from the closure Hessian at equilibrium (Proposition~\ref{prop:hessian_fisher}).
    \item \textbf{Complex.} Paper~XVIII~\citep{SmartXVIII} supplies only the amplitude-phase bookkeeping interpretation of multiplication by $i$ via the Nelson pairing $\Psi = \sqrt{\rho}\,e^{iS/\sigma^2}$; the descent of an almost-complex structure $J$ to tangent vectors of $\mathbb{CP}^{N-1}$ is part of Paper~XXI's equilibrium-tangent projective package, not an independent XVIII result.
    \item \textbf{Symplectic.} The equilibrium tangent dynamics is Hamiltonian on $\mathbb{CP}^{N-1}$ (Paper~XXI~\citep{SmartXXI}); the corresponding symplectic form $\omega$ is the imaginary part of the Hermitian structure.
\end{enumerate}

These three data $(g, J, \omega)$ are compatible (i.e., $g(v,w) = \omega(v, Jw)$) if and only if the Schr\"odinger-type equilibrium dynamics is unitary with respect to the Hermitian form $h = g + i\omega$. This is the structure assembled from Papers~XVIII and~XXI within Paper~XXI's explicit invariant-subspace hypothesis: at the equilibrium tangent, the induced dynamics is unitary on $\Hcal_\Ocal$, so $(g, J, \omega)$ assemble into a K\"ahler structure.

\begin{proposition}[Imported Fubini--Study structure with a conditional local closure-Hessian scale candidate]\label{prop:fs_assembly}
\emph{Imported inputs.}
\begin{enumerate}
    \item[\textbf{(K1)}] \emph{From Paper~XXI}, as imported standard geometric quantum mechanics on an explicitly hypothesised finite-dimensional invariant subspace $\Hcal_\Ocal$: under Paper~XXI's sector-subspace-invariance and restricted-self-adjointness hypothesis, $\mathbb{CP}^{N-1}$ carries a $U(N)$-invariant K\"ahler structure $(g_{\mathrm{XXI}}, J, \omega)$ that coincides with the standard Fubini--Study K\"ahler structure up to scale. Paper~XXI presents this as imported standard geometric QM applied to the assumed invariant subspace; its projective/K\"ahler content is not closure-derived.
    \item[\textbf{(K2)}] \emph{From Paper~XVIII}: the almost-complex structure $J$ on tangent vectors of the Nelson wavefunction is multiplication-by-$i$. Paper~XVIII itself (Section on Nelson geometric structure) states that this is a bookkeeping observation with no geometric content beyond the standard complex structure of $L^2$.
\end{enumerate}
\emph{Closure-side local computation.}
\begin{enumerate}
    \item[\textbf{(K3)}] Proposition~\ref{prop:hessian_fisher}: at a non-degenerate configuration-space equilibrium $\phi^* \in \mathbb{R}^\Ecal$, the closure Hessian $H[\phi^*]$ and the Fisher information metric $g^{\mathrm{Fisher}}$ of the location-parameter family on $\mathbb{R}^\Ecal$ satisfy $g^{\mathrm{Fisher}}(0) = (2/\sigma^2) H[\phi^*] + O(\sigma^0)$ in the Laplace regime.
\end{enumerate}
\emph{Conclusion (conditional).} Within the scope of (K1)--(K3), the standard $U(N)$-invariant K\"ahler structure of (K1) is the one used by Theorem~\ref{thm:born_unique}. Any global normalisation of that structure is conditional on a further bridge theorem identifying the closure-Hessian Fisher metric of (K3) with the real-slice Fubini--Study metric of (K1); that bridge theorem is not proved here. What is proved here is only the local, Laplace-regime, configuration-space Fisher computation of (K3), which supplies a \emph{candidate} Riemannian normalisation at a single equilibrium point $[\Psi^*]$, \emph{conditional} on such a bridge.
\end{proposition}

\begin{proof}
The $U(N)$-invariant K\"ahler structure on $\mathbb{CP}^{N-1}$ is supplied globally by (K1) from Paper~XXI, as imported standard geometric quantum mechanics under Paper~XXI's hypotheses. By Theorem~\ref{thm:chentsov}(2), any $U(N)$-invariant K\"ahler structure on $\mathbb{CP}^{N-1}$ is Fubini--Study up to a positive scalar. Paper~XXI's construction is therefore Fubini--Study up to scale. That identification does not require (K3).

The remaining question, which is the closure-side content of this paper, is whether the scale can be fixed by a closure computation. Proposition~\ref{prop:hessian_fisher} gives, via a Laplace-regime Gaussian computation in configuration space, an identification of the closure Hessian at $\phi^*$ with the Fisher information metric of the location-parameter translation family on $\mathbb{R}^\Ecal$, to leading order in $\sigma$. This is a metric on a $d$-parameter translation family in $\mathbb{R}^\Ecal$, where $d = \dim(\mathbb{R}^\Ecal)$. To turn it into a scale for the real-slice Fubini--Study metric at $[\Psi^*] \in \mathbb{CP}^{N-1}$, one would need:
\begin{itemize}
    \item an explicit diffeomorphism between the $d$-parameter translation family in configuration space and a local chart on the $(N-1)$-dimensional real slice of $\mathbb{CP}^{N-1}$ (in particular, a dimension match $d = N-1$);
    \item a separate theorem identifying the resulting configuration-space Fisher metric with the pullback of the simplex Fisher--Rao metric under the modulus-squared map $\mu_N: \mathbb{RP}^{N-1}_+ \to \Delta^{N-1}$ (Theorem~\ref{thm:chentsov}(3) supplies the simplex--$\mathbb{CP}^{N-1}$ side of this, but the configuration-space--simplex side is \emph{not} supplied by Paper~XVIII or any other cited source).
\end{itemize}
Neither step is carried out here. Consequently the scale-fixing claim is conditional on these two bridging steps, which we flag as open in Remark~\ref{rem:global_gap}. \qed
\end{proof}

\begin{remark}[What is and is not established here]\label{rem:kahler_scope}
The $U(N)$-invariant K\"ahler structure on $\mathbb{CP}^{N-1}$ is imported from Paper~XXI as standard geometric QM on the assumed invariant subspace $\Hcal_\Ocal$; it is \emph{not} derived in the present paper from the closure Hessian alone, nor is it a closure-derived output of Paper~XXI (Paper~XXI explicitly presents it as imported standard geometric QM). Proposition~\ref{prop:hessian_fisher} gives a local, leading-order identification of the closure Hessian with a configuration-space Fisher metric at the equilibrium point, but it does \emph{not} by itself fix a scale for the Fubini--Study metric on $\mathbb{CP}^{N-1}$: the pullback chain from configuration-space Fisher to real-slice Fubini--Study is not established here.
\end{remark}

\begin{remark}[Scope limitation flagged]\label{rem:global_gap}
The honest summary is: (i) the existence of a $U(N)$-invariant K\"ahler structure on $\mathbb{CP}^{N-1}$ is imported from standard geometric QM via Paper~XXI, conditional on Paper~XXI's subspace-invariance hypothesis; (ii) the closure-side content of Proposition~\ref{prop:fs_assembly} is limited to the Laplace-regime configuration-space Fisher computation of (K3); (iii) the step from (iii) to a normalisation of (i) requires further bridging not carried out here. Route~E therefore presently \emph{assembles} the geometric inputs on which a Gleason-style argument can run, rather than deriving the full $\mathbb{CP}^{N-1}$ K\"ahler structure from closure primitives.
\end{remark}

\subsection{Non-contextuality from closure locality}\label{subsec:noncontextuality}

This subsection addresses axiom (U2) of Theorem~\ref{thm:born_unique}: non-contextuality of the probability assignment. Within the VFD framework, a \emph{restricted, state-specific} form of (U2) is derivable inside the complete localised-probe-observer variant from two results in Paper~XXIX~\citep{SmartXXIX}: (i) \emph{basin locality} of closure-flow attractors, and (ii) the \emph{local measurement consistency} (LMC) theorem for the localised-probe-observer variant. The global (U2) used in Theorem~\ref{thm:born_unique} quantifies over every orthonormal basis and every $\Psi \in \mathbb{CP}^{N-1}$; upgrading the restricted form to the global one requires additionally universal completeness $\Mcl(\Ocal) = \mathbb{CP}^{N-1}$ and the assumption that every orthonormal basis is realised by a complete localised-probe observer. We summarise the restricted derivation here; the upstream proofs are in Paper~XXIX.

\begin{lemma}[Basin locality]\label{lem:basin_locality}
For a closure functional $\Fcal$ with a non-degenerate local minimum at $\phi$, the closure basin $B_\phi$ defined by the gradient flow $\dot\Phi = -\nabla\Fcal(\Phi)$ converging to $\phi$ is determined by the pair $(\Fcal, \phi)$ alone. No other minimum of $\Fcal$ enters.
\end{lemma}

This is exactly Paper~XXIX's Lemma ``Basin locality''~\citep{SmartXXIX}; we restate it here because the downstream non-contextuality proof (Theorem~\ref{thm:noncontextuality}) uses it. The content is elementary: the gradient-flow ODE involves only $\Fcal$, and the basin-membership predicate tests convergence to $\phi$ without referencing any other minimum.

Paper~XXIX introduces the \emph{localised probe observer} measurement-adapted variant: observers whose coherence constraint decomposes additively over probe subregions with off-basin silence. For this variant, Paper~XXIX proves the following theorem:

\begin{definition}[Localised probe observer and LMC; restated from Paper~XXIX]\label{def:lmc}
A \emph{localised probe observer} distinguishes a set of closure basins via a disjoint union of probe subregions $O_{\phi_j} \subset \Ecal$, with additive constraint $\Ccal_\Ocal = \sum_j c_{\phi_j}(\Phi|_{O_{\phi_j}}, \partial_{O_{\phi_j}}\Phi)$ and an off-basin-silence condition on each probe functional. Paper~XXIX proves \emph{local measurement consistency} (LMC) for this class: if two localised probe observers share a probe $(O_\phi, c_\phi)$, their constraints agree on every $\Phi \in B_\phi$, and consequently their admissible subsets of $B_\phi$ coincide.
\end{definition}

LMC is a theorem within the localised-probe-observer variant, not an axiom imposed on top of it. The present paper adopts the canonical-probe-assignment convention: observers distinguishing outcome $\phi$ use the canonical probe $(O_\phi, c_\phi)$.

\begin{theorem}[Non-contextuality of sector probability for observers complete-for-$\Psi$]\label{thm:noncontextuality}
Let $\mathcal{B} = \{\phi_1, \ldots, \phi_N\}$ and $\mathcal{B}' = \{\phi_1, \phi'_2, \ldots, \phi'_N\}$ be two orthonormal bases of $\Hcal_\Ocal$ sharing the vector $\phi_1$, corresponding to localised probe observers $\Ocal_{\mathcal{B}}, \Ocal_{\mathcal{B}'}$ (Definition~\ref{def:lmc}) that share the canonical probe for $\phi_1$. Fix a state $\Psi \in \Mcl(\Ocal_\mathcal{B}) \cap \Mcl(\Ocal_{\mathcal{B}'})$; by Definition~\ref{def:closure_manifold}, this is equivalent to requiring both observers to be complete for $\Psi$ in the sense of Remark~\ref{rem:complete_partial}, so $\int_{\Phi_{\mathrm{adm}}(\Ocal_\mathcal{B})}|\Psi|^2 \, d\mu = \int_{\Phi_{\mathrm{adm}}(\Ocal_{\mathcal{B}'})}|\Psi|^2 \, d\mu = \|\Psi\|^2$. Then for this particular $\Psi$ the normalised sector probability assignment~\eqref{eq:prob_via_integration} takes the same value in both contexts:
\begin{equation}\label{eq:noncontextual}
    P(\Pi_1, \Psi; \Ocal_{\mathcal{B}}) \;=\; P(\Pi_1, \Psi; \Ocal_{\mathcal{B}'}).
\end{equation}
The conclusion is state-specific: the equality holds for any $\Psi$ satisfying the state-completeness hypothesis with respect to both observers. To extend this to a universal non-contextuality statement for all $\Psi \in \mathbb{CP}^{N-1}$, one needs the strictly stronger \emph{universal completeness} hypothesis $\Mcl(\Ocal) = \mathbb{CP}^{N-1}$ on the relevant observer class (Remark~\ref{rem:full_CP_hypothesis}).
\end{theorem}

\begin{proof}
By Lemma~\ref{lem:basin_locality}, $B_{\phi_1}$ is the same geometric object in both contexts. By Paper~XXIX's LMC theorem applied to the canonical probe $(O_{\phi_1}, c_{\phi_1})$ shared by $\Ocal_{\mathcal{B}}$ and $\Ocal_{\mathcal{B}'}$, the observer-admissible subset of $B_{\phi_1}$ is the same in both contexts, so the numerators of~\eqref{eq:prob_via_integration} agree. By completeness, both denominators equal $\|\Psi\|^2$. Therefore the ratios agree.
\end{proof}

\begin{remark}[For partial observers, non-contextuality fails cleanly]\label{rem:partial_noncontextual}
If the two observers are \emph{partial} for $\Psi$ and have different admissibility sets (e.g.\ $\Ocal_\mathcal{B}$ admits basins $\{\Phi_1, \Phi_2\}$ while $\Ocal_{\mathcal{B}'}$ admits $\{\Phi_1, \Phi'_2\}$ with $\Phi_2 \neq \Phi'_2$), the numerators of~\eqref{eq:prob_via_integration} still agree on the shared basin $\Phi_1$ (by basin locality and LMC), but the denominators can differ. The conditional probabilities $P(\Pi_1, \Psi; \Ocal_{\mathcal{B}})$ and $P(\Pi_1, \Psi; \Ocal_{\mathcal{B}'})$ are then context-dependent. This is not a failure of the non-contextuality proof but a feature of the conditional-probability definition: probabilities conditional on different events differ, by construction. The standard-QM statement ``non-contextuality holds'' is recovered only in the complete-observer regime.
\end{remark}

\begin{corollary}[Restricted (U2) for complete, localised-probe observers, with Route~C hypotheses]\label{cor:U2_derived}
Let $P(\Pi, \Psi; \Ocal)$ denote the sector-probability assignment~\eqref{eq:prob_via_integration} realised by an observer in Paper~XXIX's localised-probe-observer variant under the canonical-probe-assignment convention. Assume further:
\begin{enumerate}
    \item[\emph{(i)}] the observer is complete for $\Psi$ (Remark~\ref{rem:complete_partial});
    \item[\emph{(ii)}] the basins $\Phi_i(\Ocal)$ are disjoint (Paper~XXIX's ambient basin-partition remark) \emph{and} jointly cover $\Phi_{\mathrm{adm}}(\Ocal)$ up to $\rho_\Psi$-null sets; the cover clause is a hypothesis of the present paper, not proved in Paper~XXIX;
    \item[\emph{(iii)}] $\Psi$ admits a sector-supported decomposition $\Psi = \sum_i c_i \phi_i$ with $|\phi_i|^2$ supported on $\Phi_i(\Ocal)$ (Route~C, Proposition~\ref{prop:born}), so in particular $\rho_\Psi$ is supported on $\bigcup_i \Phi_i(\Ocal)$.
\end{enumerate}
Then the observer-indexed family $P(\cdot, \cdot; \Ocal)$ supplies the numerator-equality and normalisation ingredients of axiom (U2) of Theorem~\ref{thm:born_unique} on the restricted class of measurement contexts described above. Obtaining the single basis-independent frame function $P(\Pi, \Psi)$ required by Theorem~\ref{thm:born_unique} from this observer-indexed family still requires the basis-realisability construction of Remark~\ref{rem:basis_realisability} together with universal completeness; the corollary by itself does not collapse the observer index.
\end{corollary}

\begin{proof}
Theorem~\ref{thm:noncontextuality} gives the non-contextuality half of (U2) under assumption (i). For the normalisation half: under (ii) and (iii), the basins $\Phi_i(\Ocal)$ partition the support of $\rho_\Psi = |\Psi|^2$, so $\sum_i \int_{\Phi_i \cap \Phi_{\mathrm{adm}}(\Ocal)}|\Psi|^2 d\mu = \int_{\Phi_{\mathrm{adm}}(\Ocal)}|\Psi|^2 d\mu$. Under (i), this denominator equals $\|\Psi\|^2$. Therefore $\sum_i P(\Pi_i, \Psi; \Ocal) = 1$ by the definition~\eqref{eq:prob_via_integration}.
\end{proof}

\begin{remark}[Scope of the reduction]\label{rem:U2_scope}
Corollary~\ref{cor:U2_derived} reduces non-contextuality from an independent axiom on the probability function to a structural condition on the class of admissible observers. This reduction is only as strong as the claim that physical measurements fall within the complete localised-probe-observer variant. If holistic or non-decomposable observers are admitted, or if the observer is partial (Remark~\ref{rem:partial}), (U2) on that enlarged class is not covered by Theorem~\ref{thm:noncontextuality}. For partial observers, VFD predicts a different, context-dependent conditional-probability structure (Remark~\ref{rem:vs_standard_qm}) --- consistent with, but extending beyond, standard QM.
\end{remark}

\subsubsection*{What this closes and what remains open}

\begin{remark}[Reduction of external axioms]\label{rem:reduction}
With Theorem~\ref{thm:noncontextuality} in place, the external inputs to Route~E are:
\begin{enumerate}
    \item The Laplace regime (Proposition~\ref{prop:hessian_fisher}).
    \item The Nelson complex structure (Paper~XVIII, Section on the geometric structure of the Nelson wavefunction).
    \item The Schr\"odinger symplectic structure and equilibrium-tangent $U(N)$ covariance (Paper~XXI, Section on the geometric structure at the equilibrium tangent).
    \item The localised-probe-observer class and canonical-probe-assignment convention (Paper~XXIX, Section on localised probe observers). LMC inside this class is proved in Paper~XXIX, not assumed.
    \item Gleason's frame-function theorem (a standard functional-analytic result).
\end{enumerate}
Non-contextuality of probability is derived only in a restricted, state-specific sense within the complete localised-probe-observer variant (Theorem~\ref{thm:noncontextuality}, Corollary~\ref{cor:U2_derived}). Upgrading this restricted statement to the global axiom (U2) used in Theorem~\ref{thm:born_unique} --- which quantifies over every orthonormal basis and every $\Psi \in \mathbb{CP}^{N-1}$ --- requires the additional assumption that every orthonormal basis is realised by a complete localised-probe observer, together with universal completeness $\Mcl(\Ocal) = \mathbb{CP}^{N-1}$. Global (U2) remains, in that sense, an external hypothesis of the forcing theorem.
\end{remark}

\begin{remark}[From a family of basis-specific observers to a single global frame function]\label{rem:basis_realisability}
Definition~\ref{def:closure_manifold} defines $\Mcl(\Ocal)$ for a single observer $\Ocal$, and Theorem~\ref{thm:noncontextuality} compares two observers $\Ocal_\mathcal{B}, \Ocal_{\mathcal{B}'}$ sharing a canonical probe. Theorem~\ref{thm:born_unique}, by contrast, uses a single basis-independent probability assignment $P(\Pi, \Psi)$ defined on all rank-one projectors $\Pi$ and all $\Psi \in \mathbb{CP}^{N-1}$, and requires (U2) to hold for every orthonormal basis simultaneously. The bridge is the following construction: given the \emph{basis-realisability} assumption that every orthonormal basis $\mathcal{B}$ of $\Hcal_\Ocal$ is realised by some complete localised-probe observer $\Ocal_\mathcal{B}$ with canonical-probe-assignment convention, define
\begin{equation}
    P(\Pi, \Psi) \;:=\; P(\Pi_\phi, \Psi; \Ocal_\mathcal{B})
\end{equation}
for any orthonormal basis $\mathcal{B}$ containing the eigenvector $\phi$ of $\Pi$. Theorem~\ref{thm:noncontextuality} (combined with universal completeness so that every $\Psi$ is state-complete for every such observer) guarantees that this value does not depend on the choice of $\mathcal{B}$, so the definition is well-posed; (U2) then follows on every basis. The basis-realisability assumption is not derived from Paper~XXIX and is not a consequence of the localised-probe-observer definition alone; it is a genuinely additional hypothesis required to turn a family of basis-specific observers into one basis-independent frame function.
\end{remark}

\begin{remark}[What remains]\label{rem:lmc_remaining}
With LMC carried as a theorem in Paper~XXIX, the main remaining question is whether every physical measurement apparatus realises a localised probe observer in the sense of Paper~XXIX's localised-probe-observer definition --- whether the probe decomposition, off-basin silence, and canonical-probe-assignment hypotheses hold for generic physical measurements. This is a question about the scope of physically realisable measurements within the VFD framework, identified as a future task (Section~\ref{sec:future}).
\end{remark}

\subsection{The forcing theorem}\label{subsec:forcing}

\begin{theorem}[Born weighting uniqueness under imported and hypothesised inputs, $N \geq 3$]\label{thm:born_unique}
Suppose $\Mcl(\Ocal) = \mathbb{CP}^{N-1}$ with $N \geq 3$ (the universal-completeness hypothesis: the observer admits every pure state as a potential closure-admissible state), equipped with the standard $U(N)$-invariant Fubini--Study K\"ahler structure imported from Paper~XXI as discussed in Proposition~\ref{prop:fs_assembly}. Let $P: \{\text{rank-1 projectors on } \Hcal_\Ocal\} \times \mathbb{CP}^{N-1} \to [0,1]$ be a probability assignment satisfying:
\begin{enumerate}
    \item[\textbf{(U1)}] \textbf{Full $U(N)$-covariance of $P$.} For every unitary $U \in U(\Hcal_\Ocal) = U(N)$,
    \begin{equation}\label{eq:U1}
        P(U\Pi U^\dagger,\, U\Psi) = P(\Pi, \Psi).
    \end{equation}
    \emph{(Paper~XXI supplies the imported equilibrium-tangent geometric $U(N)$ action on $\mathbb{CP}^{N-1}$; covariance of the probability assignment $P$ under that action is an additional hypothesis on $P$ itself, assumed here and not derived. This is the full $U(N)$ action, not the narrower sector-preserving subgroup of Proposition~\ref{prop:suff}.)}
    \item[\textbf{(U2)}] \textbf{Non-contextuality (frame-function condition).} The value $P(\Pi, \Psi)$ depends only on the projector $\Pi$ and the state $\Psi$, not on any orthogonal completion of $\Pi$ to a full resolution of the identity. Consequently, for every orthonormal basis $\{\phi_i\}_{i=1}^{N}$ of $\Hcal_\Ocal$ with $\Pi_i = |\phi_i\rangle\langle\phi_i|$,
    \begin{equation}\label{eq:U2}
        \sum_{i=1}^{N} P(\Pi_i, \Psi) = 1 \qquad \forall \Psi \in \Mcl.
    \end{equation}
    \emph{(U2) is taken as a hypothesis of this theorem. A restricted, state-specific form of (U2) is derived within the Paper~XXIX localised-probe-observer variant from basin locality (Lemma~\ref{lem:basin_locality}) and local measurement consistency (Definition~\ref{def:lmc}) via Theorem~\ref{thm:noncontextuality}; that restricted form does not by itself suffice for the Gleason step, which requires (U2) to hold globally over all orthonormal bases and all $\Psi \in \mathbb{CP}^{N-1}$. Upgrading the restricted form to (U2) requires either adopting (U2) outright, or combining Theorem~\ref{thm:noncontextuality} with \emph{both} universal completeness $\Mcl(\Ocal) = \mathbb{CP}^{N-1}$ \emph{and} the assumption that every orthonormal basis is realised by a complete localised-probe observer. See Corollary~\ref{cor:U2_derived} and Remark~\ref{rem:full_CP_hypothesis}.)}
    \item[\textbf{(U3)}] \textbf{Smoothness.} $\Psi \mapsto P(\Pi, \Psi)$ is continuous on $\mathbb{CP}^{N-1}$ for each fixed $\Pi$.
    \item[\textbf{(U4)}] \textbf{Eigenstate certainty.} If $\Pi = |\phi\rangle\langle\phi|$ and $\Psi = \phi$ (i.e., $\Pi\Psi = \Psi$), then $P(\Pi, \Psi) = 1$.
\end{enumerate}
Then $P$ is the Born measure:
\begin{equation}\label{eq:born_forced}
    P(\Pi, \Psi) \;=\; \frac{\langle\Psi | \Pi | \Psi\rangle}{\langle\Psi|\Psi\rangle}.
\end{equation}
\end{theorem}

\begin{remark}[Why $\Mcl(\Ocal) = \mathbb{CP}^{N-1}$ is a genuine hypothesis, not a simplification]\label{rem:full_CP_hypothesis}
Definition~\ref{def:closure_manifold} defines $\Mcl(\Ocal)$ as a \emph{subset} of $\mathbb{CP}^{N-1}$: the projective states whose Nelson density is supported on $\Phi_{\mathrm{adm}}(\Ocal)$. The theorem hypothesis strengthens this to equality: every pure projective state is observer-admissible. This is load-bearing. Without equality: (a) the $U(N)$-orbit argument in Step~1 below does not reach states outside $\Mcl$, so one cannot reduce $P(\Pi, \Psi)$ to a scalar function of the full range of overlaps $\langle\Psi|\Pi|\Psi\rangle$; (b) Gleason's frame-function theorem in Step~2 requires all orthonormal bases of $\Hcal_\Ocal$ to be within the theorem's scope.

Whether $\Mcl = \mathbb{CP}^{N-1}$ holds is a question about the observer's coherence constraint $\Ccal_\Ocal$, not about the closure dynamics alone: it requires every pure superposition $\Psi = \sum_i c_i \phi_i$ to have Nelson density $\rho_\Psi$ with $\mathrm{supp}(\rho_\Psi) \subseteq \Phi_{\mathrm{adm}}(\Ocal)$. This is an observer-side condition: the observer must admit all sector superpositions as coherent incoming states. The present theorem takes this condition as a \emph{hypothesis}; it is not derived here. In particular, Paper~XXI's equilibrium-tangent unitarity gives a $U(N)$ action on $\mathbb{CP}^{N-1}$ but does \emph{not} by itself establish the support condition on $\rho_\Psi$: a generic superposition may have density escaping the observer-admissible set. Whether physical measurement observers satisfy $\Mcl = \mathbb{CP}^{N-1}$ is an open question identified in Section~\ref{sec:future}; the theorem is conditional on it.
\end{remark}

\begin{proof}
\textbf{Step 1: Reduction to a function of the projector overlap.} By (U1) and (U2), the scalar $P(\Pi, \Psi)$ is a \emph{frame function} in Gleason's sense: its value depends on the single projector $\Pi$ and the state $\Psi$, not on the orthogonal completion. By (U1), $P$ depends only on the $U(N)$-orbit of the pair $(\Pi, \Psi)$, which for rank-1 $\Pi$ and unit $\Psi$ is parameterised by the single invariant $x = \langle\Psi | \Pi | \Psi\rangle \in [0,1]$. Therefore there exists a continuous function $g: [0,1] \to [0,1]$ with
\begin{equation}
    P(\Pi, \Psi) = g(x), \qquad x = \langle\Psi|\Pi|\Psi\rangle.
\end{equation}

\textbf{Step 2: Gleason plus the frame-sum constraint force $g$ affine.} Gleason's frame-function theorem for continuous frame functions on a Hilbert space of dimension $\geq 3$ says: if $f$ is a non-negative continuous function on the unit sphere of $\Hcal$ with the property that $\sum_i f(\phi_i)$ takes the same value for every orthonormal basis $\{\phi_i\}$, then $f(\Psi) = \langle\Psi|A|\Psi\rangle$ for some self-adjoint operator $A$~\citep{Gleason1957,Peres1995,CookKeaneMoran1985}. For any fixed unit $\Psi$, apply this with $f_\Psi(\phi) := g(|\langle\phi|\Psi\rangle|^2) = g(\langle\Psi|\Pi_\phi|\Psi\rangle)$, where $\Pi_\phi = |\phi\rangle\langle\phi|$; by (U2), $\sum_i f_\Psi(\phi_i) = 1$ for every orthonormal basis, so $f_\Psi$ is a frame function, and by Gleason's theorem $f_\Psi(\phi) = \langle\phi|A_\Psi|\phi\rangle$ for some self-adjoint $A_\Psi$. Since $f_\Psi(\phi)$ depends on $\phi$ only through the scalar $|\langle\phi|\Psi\rangle|^2$, the operator $A_\Psi$ must commute with every rotation fixing $|\Psi\rangle$, which forces $A_\Psi = a_\Psi \Pi_\Psi + b_\Psi (I - \Pi_\Psi)$ for scalars $a_\Psi, b_\Psi$. Comparing with $g(x) = \langle\phi|A_\Psi|\phi\rangle = a_\Psi x + b_\Psi(1-x) = (a_\Psi - b_\Psi) x + b_\Psi$ (where $x = |\langle\phi|\Psi\rangle|^2$) shows that $g$ is affine, $g(x) = \alpha x + \beta$ with $\alpha = a_\Psi - b_\Psi$ and $\beta = b_\Psi$; these constants are independent of $\Psi$ because $g$ was fixed at the start. The entire content of Step~2 is therefore imported from Gleason's theorem; the remainder of this proof establishes that the VFD framework satisfies Gleason's hypotheses and fixes his free constants.

\textbf{Step 3: (U4) fixes the affine constants.} Write $g(x) = \alpha x + \beta$ from Step~2. Applying (U2) to any orthonormal basis containing $\Pi$ and taking $\Psi = \phi$ (the eigenvector of $\Pi$) gives
\begin{equation}
    g(1) + (N-1) g(0) = 1, \quad \text{i.e.,} \quad (\alpha + \beta) + (N-1)\beta = \alpha + N\beta = 1.
\end{equation}
Axiom (U4) gives $g(1) = \alpha + \beta = 1$. Subtracting the two equations yields $(N-1)\beta = 0$, and since $N \geq 3$ we conclude $\beta = 0$, $\alpha = 1$. Therefore
\begin{equation}
    P(\Pi, \Psi) = \langle\Psi|\Pi|\Psi\rangle,
\end{equation}
which is the Born rule. Note that (U4) is the eigenstate-certainty half of the Born rule itself, and is the minimal additional input beyond (U1)--(U3) needed to pin down both $\alpha$ and $\beta$: without (U4), the two equations $\alpha + \beta \leq 1$ and $\alpha + N\beta = 1$ underdetermine the two constants. An equivalent alternative would be the null-state axiom $g(0) = 0$ (probability of an outcome whose eigenstate is orthogonal to the state is zero); we adopt (U4) as the more physical statement.
\end{proof}

\begin{corollary}[Exclusion of $|c|^p$ for $p \neq 2$]\label{cor:p_excluded}
Let $N \geq 3$. For $p > 0$, consider the assignment
\begin{equation}
    Q(\Pi_\phi, \Psi;\, \{\phi_j\}_{j=1}^{N}) \;=\; \frac{|\langle\phi|\Psi\rangle|^p}{\sum_{j=1}^{N} |\langle\phi_j|\Psi\rangle|^p},
\end{equation}
where $\Pi_\phi = |\phi\rangle\langle\phi|$ and $\{\phi_j\}$ is any orthonormal basis containing $\phi$. This assignment satisfies axiom (U2) of Theorem~\ref{thm:born_unique} --- that is, $Q(\Pi_\phi, \Psi;\, \cdot)$ is independent of the basis completion $\{\phi_j\}$ --- if and only if $p = 2$.
\end{corollary}

\begin{proof}
($\Leftarrow$) For $p = 2$, Parseval's identity gives $\sum_j |\langle\phi_j|\Psi\rangle|^2 = \|\Psi\|^2 = 1$ in every orthonormal basis, so the denominator is basis-independent and $Q = |\langle\phi|\Psi\rangle|^2 = \langle\Psi|\Pi_\phi|\Psi\rangle$ depends only on $(\Pi_\phi, \Psi)$.

($\Rightarrow$) For $p \neq 2$, we exhibit a concrete $N = 3$ counterexample that embeds in any $N \geq 3$. Let $\{e_1, e_2, e_3\}$ be an orthonormal basis, fix $\Pi_{e_1} = |e_1\rangle\langle e_1|$, and take $\Psi = (e_1 + e_2 + e_3)/\sqrt{3}$. Two orthonormal completions of $\{e_1\}$ on the three-dimensional subspace $\mathrm{span}\{e_1, e_2, e_3\}$ are:
\begin{itemize}
    \item \textbf{Completion~A:} $\{e_1, e_2, e_3\}$.
    \item \textbf{Completion~B:} $\{e_1, f_+, f_-\}$, where $f_\pm = (e_2 \pm e_3)/\sqrt{2}$.
\end{itemize}
The inner products with $\Psi$ are:
\begin{equation}
    \langle e_1|\Psi\rangle = \tfrac{1}{\sqrt{3}}, \quad \langle e_2|\Psi\rangle = \tfrac{1}{\sqrt{3}}, \quad \langle e_3|\Psi\rangle = \tfrac{1}{\sqrt{3}}, \quad \langle f_+|\Psi\rangle = \sqrt{\tfrac{2}{3}}, \quad \langle f_-|\Psi\rangle = 0.
\end{equation}
The denominators for $Q$ are
\begin{align}
    D_A &= 3 \cdot (1/\sqrt{3})^p = 3^{1 - p/2}, \\
    D_B &= (1/\sqrt{3})^p + (\sqrt{2/3})^p + 0 = 3^{-p/2}\bigl(1 + 2^{p/2}\bigr),
\end{align}
and the common numerator is $|\langle e_1|\Psi\rangle|^p = 3^{-p/2}$. Therefore
\begin{equation}
    Q_A = \frac{3^{-p/2}}{3^{1 - p/2}} = \frac{1}{3},
    \qquad
    Q_B = \frac{3^{-p/2}}{3^{-p/2}(1 + 2^{p/2})} = \frac{1}{1 + 2^{p/2}}.
\end{equation}
These coincide precisely when $1 + 2^{p/2} = 3$, i.e., $2^{p/2} = 2$, i.e., $p = 2$. For any $p \neq 2$, $Q_A \neq Q_B$, so $Q$ depends on the basis completion of $\Pi_{e_1}$ and violates (U2).
\end{proof}

\begin{remark}[Two-dimensional case, $N = 2$: not covered by Theorem~\ref{thm:born_unique}]\label{rem:N_two}
Gleason's theorem requires $\dim \Hcal \geq 3$ and does not apply directly at $N = 2$: in dimension two, the frame-function argument (Step~2 of Theorem~\ref{thm:born_unique}) admits continuous solutions beyond the quadratic form, so the Step~2 conclusion does not hold for $N = 2$. The statement of Theorem~\ref{thm:born_unique} is therefore scoped to $N \geq 3$ only. For $N = 2$, a separate uniqueness argument would be needed; a standard route on the Bloch sphere $\mathbb{CP}^1 \cong S^2$ is to take $P(\Pi, \Psi)$ to be (i) continuous, (ii) valued in $[0,1]$, (iii) $SU(2)$-invariant under the diagonal action, (iv) antipodally normalised ($P(\Pi, \Psi) + P(-\Pi, \Psi) = 1$ for antipodal projectors), and (v) certain on eigenstates ($P(\Pi, \phi) = 1$ when $\Psi$ is the Bloch vector of $\Pi$), in which case the solution is conjectured to be the latitude function $\tfrac12(1 + \hat{n}\cdot\hat{m}) = |\langle\phi|\Psi\rangle|^2$. We do not prove this uniqueness claim here: it requires a direct analysis of $SU(2)$-invariant continuous functions on $S^2 \times S^2$ satisfying (iv) and (v), which we have not carried out. The $N = 2$ case is therefore open within the present paper.
\end{remark}

\begin{remark}[Relationship to Gleason and to Chentsov--Petz]\label{rem:gleason_relation}
The proof of Theorem~\ref{thm:born_unique} is essentially Gleason's: the key technical content is the frame-function theorem, which VFD does not reprove. Route~E's contribution is not to reprove uniqueness but to exhibit a coherent dependency chain in which the Hilbert-space and projective/K\"ahler structures Gleason's theorem assumes are packaged together: the finite-dimensional invariant-subspace Hilbert space comes from Paper~XXI (standard geometric QM under its explicit hypothesis); the Nelson-wavefunction complex-structure bookkeeping comes from Paper~XVIII; the restricted form of non-contextuality is supplied by Paper~XXIX's LMC theorem inside the complete localised-probe-observer variant. Route~E does \emph{not} claim that these structures are produced from the full nonlinear closure dynamics; Paper~XXI and Paper~XVIII are explicit that they are not. The independent information-geometric content is Theorem~\ref{thm:chentsov}: once the imported $U(N)$-invariant K\"ahler structure is identified with Fubini--Study up to scale, the push-forward to the simplex realises the Fisher--Rao metric. Chentsov's theorem certifies this push-forward as the canonical one; it is a cross-check on the imported structure rather than an independent uniqueness driver for the Born rule.
\end{remark}

\subsection{What Route E adds --- and where external input enters}\label{subsec:route_e_delta}

Route~C establishes conditional existence: in the equilibrium tangent regime, the stationary sector measure of the closure dynamics realises the Born distribution. Route~D establishes conditional uniqueness given Hilbert structure, via Gleason's theorem. Route~E's contribution is neither purely existential nor purely structural: it assembles, layer by layer, the geometric data on which a Gleason-type uniqueness argument becomes applicable, combining closure-side and imported geometric inputs.

The Route~E chain and its external inputs are:

\begin{itemize}
    \item \textbf{Statistical parameter-space structure} (Definition~\ref{def:closure_manifold}): defined using Paper~XXIX's admissibility together with Paper~XVIII's Nelson pairing. No well-separated-sector hypothesis is needed at this definitional stage; well-separation enters only at Route~C's sector-supported decomposition.
    \item \textbf{Riemannian metric} (Proposition~\ref{prop:hessian_fisher}): derived from the closure Hessian at equilibrium \emph{in the Laplace regime}. External input: the Laplace (small-$\sigma$) approximation.
    \item \textbf{Complex structure} (Proposition~\ref{prop:fs_assembly}): imported from Paper~XVIII as a bookkeeping observation on the Nelson wavefunction. External input: Paper~XVIII's Section on the geometric structure of the Nelson wavefunction.
    \item \textbf{Symplectic structure and $U(N)$ covariance} (Proposition~\ref{prop:fs_assembly}): imported from Paper~XXI as standard geometric QM on an assumed invariant subspace. External input: Paper~XXI's Section on geometric structure at the equilibrium tangent.
    \item \textbf{Non-contextuality of probability} (axiom (U2) of Theorem~\ref{thm:born_unique}): \emph{derived only in a restricted, state-specific form} within the complete localised-probe-observer variant, from basin locality (Lemma~\ref{lem:basin_locality}) and local measurement consistency (Definition~\ref{def:lmc}), via Theorem~\ref{thm:noncontextuality}. The global (U2) used by Theorem~\ref{thm:born_unique} requires additionally universal completeness and the assumption that every orthonormal basis is realised by a complete localised-probe observer; both remain external hypotheses.
    \item \textbf{Frame-function theorem} (Step~2 of the proof): imported from Gleason's theorem~\citep{Gleason1957,Peres1995,CookKeaneMoran1985}. External input: a standard theorem of functional analysis on Hilbert space.
\end{itemize}

Route~E is therefore a construction showing that, given Papers~XVIII and~XXI (the $U(N)$-invariant K\"ahler structure on $\mathbb{CP}^{N-1}$ imported from standard geometric QM on Paper~XXI's assumed invariant subspace) and the global non-contextuality hypothesis (U2), Gleason's theorem identifies the Born rule as the unique probability assignment satisfying (U1)--(U4) for $N \geq 3$. The Chentsov--Petz content is structural, not load-bearing: it certifies that the imported $U(N)$-invariant K\"ahler structure is specifically Fubini--Study up to scale, providing an information-geometric cross-check but not itself driving the uniqueness proof. The scale-fixing bridge theorem concerns only the candidate local normalisation via the closure Hessian; it is not a condition for the Gleason uniqueness step. The non-contextuality derivation (Section~\ref{subsec:noncontextuality}) delivers a restricted, state-specific form of (U2) within the complete localised-probe-observer variant; upgrading to the global (U2) used by Theorem~\ref{thm:born_unique} requires universal completeness together with the assumption that every orthonormal basis is realised by a complete localised-probe observer.

\begin{remark}[Scope of the uniqueness]\label{rem:scope}
The uniqueness result is conditional. Theorem~\ref{thm:born_unique} runs only after universal completeness identifies $\Mcl(\Ocal)$ with the full $\mathbb{CP}^{N-1}$; the equilibrium tangent regime (Papers~XVIII--XXI, under Paper~XXI's sector-subspace invariance hypothesis) and the Laplace approximation holding locally at an equilibrium state are used only to license the imported geometric inputs on which the theorem is then applied. The theorem itself assumes global (U1)--(U4) on all of $\mathbb{CP}^{N-1}$ with $N \geq 3$; Route~C's well-separated-sector hypothesis is \emph{not} part of this uniqueness theorem, but enters only as part of the separate Route~C existence route and of the basis-realisability bridge that upgrades Paper~XXIX's restricted non-contextuality to the global (U2) hypothesis. The $N = 2$ case is not covered by this theorem and is not proved in the present paper (Remark~\ref{rem:N_two}). The result does not cover the pre-equilibrium nonlinear regime of Paper~XIX~\citep{SmartXIX}, nor does it establish that physical measurement interactions generically produce the required sector structure. Within the stated scope (including $N \geq 3$), any globally defined probability assignment satisfying (U1)--(U4) is forced to the Born form; outside that scope, including $N = 2$, the question is open.
\end{remark}

\begin{remark}[Comparison with Zurek's envariance]\label{rem:zurek}
Zurek's envariance argument~\citep{Zurek2005} derives the Born rule from the invariance of entangled states under local unitaries on an environment. The Route~E argument here is structurally similar --- both proceed from a unitary invariance --- but with a different substrate: envariance requires an entangled state with an environment; Route~E requires the imported K\"ahler structure of Paper~XXI (Proposition~\ref{prop:fs_assembly}) and, ultimately, the global non-contextuality hypothesis (U2). Paper~XXIX's LMC theorem, via Theorem~\ref{thm:noncontextuality}, supplies only a restricted, state-specific precursor to (U2); it does not by itself discharge the hypothesis used by Theorem~\ref{thm:born_unique}. The comparison is heuristic only; this paper does not derive the measurement-theoretic origin of the required global unitary covariance.
\end{remark}

%==================================================================
\section{Toy Model: Two-Basin Weighting}\label{sec:toy}
%==================================================================

This is a schematic effective model illustrating sector-mass arithmetic. It is not derived from the full 600-cell construction or from the double-well closure potential alone: the asymmetric wavefunction used below is stipulated, not produced by the closure dynamics itself. Its purpose is to illustrate the mechanism by which sector-supported amplitudes yield Born-like weights, not to demonstrate that the closure dynamics generates such amplitudes.

\begin{example}[Two-basin model]\label{ex:two_basin}
Consider a one-dimensional closure functional with two minima:
\begin{equation}
    \Fcal(\psi) = (\psi^2 - 1)^2
\end{equation}
This has minima at $\psi^* = \pm 1$ with $\Fcal(\pm 1) = 0$, and a maximum at $\psi = 0$ with $\Fcal(0) = 1$.

\textbf{Symmetric case.} The stationary distribution is $\Pst(\psi) \propto e^{-2(\psi^2-1)^2/\sigma^2}$. By symmetry, $\int_{-\infty}^0 \Pst = \int_0^\infty \Pst = 1/2$. The two basins have equal weight: $P_1 = P_2 = 1/2$.

\textbf{Asymmetric case (schematic).} We now insert, \emph{by hand}, an asymmetric Nelson wavefunction:
\begin{equation}
    \Psi(\psi) = a\, \phi_1(\psi) + b\, \phi_2(\psi),
\end{equation}
where $\phi_1$ is a normalised Gaussian localised near $\psi = -1$, $\phi_2$ is localised near $\psi = +1$, and $|a|^2 + |b|^2 = 1$. This state is \emph{not} the symmetric stationary state of the double-well $\Fcal(\psi) = (\psi^2-1)^2$; it is a stipulated asymmetric state used only to illustrate the sector-mass arithmetic when a state is sector-supported.

Granting the stipulation, the stationary-measure weights of this $\Psi$ are:
\begin{align}
    w_1 &= \int_{-\infty}^{0} |\Psi|^2\, d\psi \;\approx\; |a|^2 \quad \text{(cross terms negligible by separation)} \\
    w_2 &= \int_{0}^{\infty} |\Psi|^2\, d\psi \;\approx\; |b|^2.
\end{align}
Therefore $P_1 = |a|^2, \; P_2 = |b|^2$, which matches the Born-rule algebra for a two-outcome split.

\textbf{What this example does and does not show.} The calculation above establishes the sector-mass arithmetic under the stipulation that $\Psi$ is a sector-supported superposition with fixed amplitudes $(a,b)$. It \emph{does not} derive those amplitudes from the closure dynamics of the double-well potential; in particular, the double-well's actual stationary state is symmetric, and producing the asymmetric $\Psi$ would require either an external boundary condition or coupling to an observer constraint that breaks the symmetry. The example illustrates the mechanism by which sector-localised amplitudes yield Born-like weights; it does not constitute a derivation that the closure dynamics itself produces asymmetric states with prescribed $(a,b)$.
\end{example}

%==================================================================
\section{Relationship to Standard Quantum Theory}\label{sec:comparison}
%==================================================================

\subsection{What is preserved}

For complete observers in the equilibrium, well-separated, sector-supported, stationary-sampling regime analysed here, and under the Route~E completeness and basis-realisability assumptions, the operational predictions agree with standard quantum mechanics. The VFD proposal does not modify the Born rule in that regime; it seeks to explain why the Born rule has the form it does. Amplitudes remain the effective computational tool for complete-observer measurements. For partial observers, VFD predicts conditional probabilities that differ from the standard Born assignment by a normalisation over the admissible set (Remarks~\ref{rem:partial}, \ref{rem:vs_standard_qm}); this is an extension of, not a contradiction with, standard QM, because standard QM implicitly assumes complete observers and handles partial measurements by bolting on extra machinery (POVMs, no-click outcomes). The combined proposal is therefore: a \textit{substrate explanation} of the Born rule for complete measurements (via the conditional equilibrium-sector argument of Route~C), plus a \textit{structural prediction} of conditional probabilities for partial measurements (from the present observer-conditioning formalism). The two parts share a common framework, but are not delivered by a single closure-dynamical theorem.

\subsection{Comparison with other approaches}

\begin{table}[ht]
\centering
\caption{Born-rule status across interpretive frameworks.}
\label{tab:born}
\begin{tabular}{@{}lll@{}}
\toprule
Framework & Born rule status & Mechanism \\
\midrule
Copenhagen & Postulated axiom & None (primitive) \\
Many-worlds & Derived (decision/symmetry) & Branch counting / envariance \\
Bohmian mechanics & Derived (quantum equilibrium) & Typicality of initial conditions \\
Objective collapse & Modified / stochastic & Added stochastic law \\
VFD (this paper) & Conditional for complete observers under the stated Route~C and Route~E hypotheses; observer-conditional for partial observers & Stationary sector measure + Gleason uniqueness on imported K\"ahler/Hilbert structure, normalised over admissible set \\
\bottomrule
\end{tabular}
\end{table}

\begin{remark}
The VFD route is closest in spirit to the Bohmian quantum-equilibrium approach: both derive the Born distribution from an underlying dynamical measure. The key difference is that the Bohmian approach assumes pilot-wave dynamics and derives $|\psi|^2$ as the equivariant measure, while VFD derives $|\Psi|^2$ as the stationary measure of the closure dynamics. The Bohmian approach requires a separate justification for why the universe is in quantum equilibrium; the VFD approach identifies the Paper~XIV equilibrium as the relevant stationary regime, while generic convergence to it from arbitrary initial conditions remains open (Section~\ref{sec:discussion}'s discussion of the equilibrium assumption).
\end{remark}

%==================================================================
\section{Route Comparison and Synthesis}\label{sec:synthesis}
%==================================================================

\begin{table}[ht]
\centering
\caption{Comparison of the five weighting routes.}
\label{tab:routes}
\begin{tabular}{@{}llll@{}}
\toprule
Route & Weight $w_i$ & Born-like? & Status \\
\midrule
A: Basin volume & $\mu(B_i \cap \Phi_{\mathrm{adm}})$ & Not necessarily & Motivational prototype \\
B: Stability depth & $e^{-\beta \Fcal[\Phi_i^*]}$ & Boltzmann, not Born & Motivational prototype \\
C: Stationary measure & $\int_{\Phi_i} |\Psi|^2\, d\mu$ & Yes (at equilibrium) & Main conditional derivation \\
D: Gleason bridge & $|\langle \phi_i | \psi \rangle|^2$ & Yes (by theorem) & Structural confirmation (given Hilbert) \\
E: K\"ahler uniqueness & $|c_i|^2$ (unique) & Forced (for $N{\geq}3$) & Uniqueness programme via imported K\"ahler structure + Gleason \\
\bottomrule
\end{tabular}
\end{table}

The five routes fall into three strata, matching the three levels introduced in Section~\ref{sec:intro}:

\begin{itemize}
    \item \textbf{Motivational prototypes (A, B).} These illuminate why probability must be geometric in character. Basin volume explains why outcomes are not automatically equiprobable; stability depth explains why deeper minima are weighted more heavily. Neither produces quadratic weighting in isolation and neither is advanced as a complete candidate.
    \item \textbf{Conditional derivation (C).} In the equilibrium tangent regime, the stationary measure of the closure dynamics \emph{realises} the Born distribution: $|\Psi|^2 = \Pst$, and sector probabilities are inherited as $|c_i|^2$ under the stated assumptions. This is the main existence result.
    \item \textbf{Uniqueness arguments (D, E).} Route~D gives Gleason-style uniqueness once Hilbert structure is granted. Route~E assembles a geometric dependency chain: the $U(N)$-invariant K\"ahler structure on $\mathbb{CP}^{N-1}$ is imported from Paper~XXI as standard geometric QM on the assumed invariant subspace; the complex structure is the bookkeeping observation from Paper~XVIII; Petz's theorem identifies the imported structure as Fubini--Study up to scale; and Gleason's frame-function theorem for $N \geq 3$ fixes the unique probability assignment as the Born measure under (U1)--(U4) and universal completeness. The closure Hessian (Proposition~\ref{prop:hessian_fisher}) supplies only a candidate local scale, conditional on an unproved bridge theorem from configuration-space Fisher to real-slice Fubini--Study; it does not play an operative role in the proof of Theorem~\ref{thm:born_unique}, which is Gleason's. The $N = 2$ case is not covered here (Remark~\ref{rem:N_two}). The conditional nature of this uniqueness --- what is imported vs what is hypothesised vs what is proved here --- is made explicit in Section~\ref{subsec:route_e_delta}.
\end{itemize}

The combined Routes~C and~E argument establishes the Born rule conditionally within the equilibrium closure regime: existence from the stationary measure (plus the sampling hypothesis of Proposition~\ref{prop:born} assumption~(4) and the state-completeness hypothesis of assumption~(5)); uniqueness from Gleason's theorem applied under axioms (U1)--(U4) and universal completeness, for observers that satisfy $\Mcl(\Ocal) = \mathbb{CP}^{N-1}$ and fall within Paper~XXIX's localised-probe-observer variant. The imported $U(N)$-invariant K\"ahler structure of Paper~XXI and the local closure-Hessian computation of Proposition~\ref{prop:hessian_fisher} package and interpret the imported geometry but do not enter the uniqueness proof itself (Remark~\ref{rem:global_gap}). Route~D provides an independent structural check and a bridge to the orthodox-quantum literature. A restricted, state-specific form of non-contextuality is derived within Paper~XXIX's localised-probe-observer variant (Corollary~\ref{cor:U2_derived}, citing Paper~XXIX); the global (U2) used by Theorem~\ref{thm:born_unique} is obtained from that restricted form only in combination with universal completeness $\Mcl = \mathbb{CP}^{N-1}$ and the basis-realisability assumption that every orthonormal basis of $\Hcal_\Ocal$ is realised by some complete localised-probe observer (Remark~\ref{rem:basis_realisability}). The remaining observer-structural inputs are (a) the stipulation that physical measurements fall within the localised-probe-observer variant, (b) universal completeness, and (c) basis-realisability. The Route~C sampling hypothesis is a separate measurement-theoretic external assumption, not an observer-structural one. Each is a substantive external assumption, not a derived consequence of the present paper.

%==================================================================
\section{Discussion and Limitations}\label{sec:discussion}
%==================================================================

\subsection{What is established}

\begin{itemize}
    \item Admissibility alone does not determine outcome probabilities; a weighting measure is required (Section~\ref{sec:gap}).
    \item Five weighting routes are examined, with explicit constructions for each. Routes~A and~B are motivational prototypes; Route~C is the main conditional derivation; Route~D is a structural confirmation; Route~E assembles imported geometric inputs and applies Gleason's frame-function theorem, with information-geometric content serving as structural cross-check.
    \item \textbf{Conditional existence (Route~C).} Under the following five hypotheses --- equilibrium is reached, sectors are well-separated, the state admits a sector-supported decomposition, measurement samples the constrained stationary ensemble, and the observer is state-complete for the state in question (Nelson density supported in the admissible set) --- the stationary-measure route yields the algebraic identity $P_i = |c_i|^2$ up to sector-overlap corrections (Proposition~\ref{prop:born}). The sampling hypothesis and the state-completeness hypothesis are both identified explicitly in Section~\ref{sec:stationary} as working assumptions, not inherited consequences; dropping state-completeness leaves the observer-conditional form $|c_i|^2/\sum_{j\in\mathrm{adm}}|c_j|^2$.
    \item \textbf{Gleason uniqueness on imported geometric packaging (Route~E), in three steps.}
    \begin{enumerate}
        \item The closure Hessian at equilibrium \emph{locally} realises the Fisher information metric in the Laplace approximation (Proposition~\ref{prop:hessian_fisher}, with leading-order identification and $O(\sigma^0)$ corrections).
        \item The $U(N)$-invariant K\"ahler structure on $\Mcl = \mathbb{CP}^{N-1}$ is imported from Paper~XXI as standard geometric QM on the assumed invariant subspace, and by Petz's uniqueness theorem coincides with Fubini--Study up to scale; the local Fisher computation of item~1 supplies a candidate scale at the equilibrium point, conditional on a bridge theorem not proved in this paper (Proposition~\ref{prop:fs_assembly}, Remark~\ref{rem:global_gap}).
        \item On this K\"ahler manifold with $N \geq 3$, any probability assignment satisfying (U1)--(U4) --- unitary covariance, non-contextuality, continuity, and eigenstate certainty --- is the Born measure, by Gleason's frame-function theorem (Theorem~\ref{thm:born_unique}).
    \end{enumerate}
    \item \textbf{Exclusion of $|c|^p$ for $p \neq 2$.} The alternative $Q_i = |c_i|^p/\sum_j |c_j|^p$ is non-contextual only for $p = 2$, by Parseval's identity (Corollary~\ref{cor:p_excluded}).
    \item \textbf{$N = 2$ not covered here.} The two-dimensional case falls outside Gleason's hypotheses and is not proved in the present paper; a direct $SU(2)$-invariant Bloch-sphere argument is sketched in Remark~\ref{rem:N_two} as a candidate route, but the uniqueness claim it would yield is not verified here.
    \item \textbf{Structural redundancy (Route~D).} The Gleason bridge provides an independent confirmation of the uniqueness half, using Hilbert structure directly as input rather than constructing it from closure.
    \item The sector-mass arithmetic is illustrated in a two-basin toy model (Example~\ref{ex:two_basin}), with the asymmetric state stipulated by hand rather than derived from the closure dynamics.
\end{itemize}

\subsection{What is imported from prior papers}

The uniqueness argument in Route~E relies on the following prior-paper results. These are flagged explicitly so the scope of the derivation is transparent:

\begin{itemize}
    \item \textbf{Laplace regime} (Proposition~\ref{prop:hessian_fisher}): the closure Hessian identification with the Fisher metric is a leading-order statement in the small-$\sigma$ Gaussian approximation, not a global equality.
    \item \textbf{Nelson complex-structure bookkeeping} (Paper~XVIII, Section on geometric structure of the Nelson wavefunction): the multiplication-by-$i$ on amplitude-phase pairs yields the almost-complex structure $J$, presented in Paper~XVIII as a bookkeeping observation with no independent geometric content beyond the standard complex structure of $L^2$.
    \item \textbf{Schr\"odinger symplectic structure and $U(N)$ covariance} (Paper~XXI, Section on geometric structure at the equilibrium tangent): the linearised Schr\"odinger dynamics is a Hamiltonian flow on $\mathbb{CP}^{N-1}$ with $U(N)$ acting as K\"ahler isometries. Paper~XXI presents this as standard geometric QM imported under its explicit sector-subspace invariance and restricted-self-adjointness hypothesis; it is not a closure-derived result.
    \item \textbf{Localised probe observers and LMC} (Paper~XXIX, Section on localised probe observers): observers with additive basin-local coherence constraints satisfying off-basin silence. LMC is a theorem within this class, proved in Paper~XXIX.
    \item \textbf{Gleason's frame-function theorem} (Step~2 of the forcing-theorem proof): a standard result of functional analysis.
\end{itemize}

Non-contextuality of probability --- previously the single Gleason-style axiom that Route~E imported --- is derived only in a restricted, state-specific sense within the complete localised-probe-observer variant (Theorem~\ref{thm:noncontextuality}, Corollary~\ref{cor:U2_derived}). The global (U2) axiom used in Theorem~\ref{thm:born_unique}, which quantifies over every orthonormal basis and every state in $\mathbb{CP}^{N-1}$, is \emph{not} delivered by this restricted derivation alone; upgrading to the global form requires both universal completeness $\Mcl(\Ocal) = \mathbb{CP}^{N-1}$ and the assumption that every orthonormal basis is realised by a complete localised-probe observer. What remains on the observer side is therefore (a) the structural condition that physical observers are localised probe observers, and (b) the basis-realisability and completeness assumptions needed to lift the restricted non-contextuality to the global (U2).

\subsection{What is not established}

\begin{itemize}
    \item Proof that every physical measurement apparatus realises a localised probe observer in the sense of Paper~XXIX's localised-probe-observer definition. With LMC carried upstream as a theorem for that variant, the remaining observer-side question is whether the probe decomposition and off-basin silence hypotheses hold for generic physical measurements. A holistic (non-decomposable) observer would fall outside the variant and be uncovered by the derivation.
    \item Proof that physical measurement interactions generically produce the well-separated sector structure assumed in Route~C and used by Route~E. This is a decoherence question.
    \item Extension of the uniqueness argument to continuous spectra or infinite-dimensional configuration spaces. The Chentsov and Gleason results have extensions to this regime~\citep{Petz1996}, but their application to VFD has not been spelled out here.
    \item A concrete construction of the sector decomposition $\Psi = \sum_i c_i \phi_i$ from 600-cell geometry.
    \item Proof that the system reaches equilibrium on physically relevant timescales. (Analogous to the quantum-equilibrium hypothesis of Bohmian mechanics.)
    \item The relationship between the coherence threshold $\epsilon$ and the basin structure.
    \item Corrections to the Born rule away from equilibrium. The nonlinear residual of Paper~XIX~\citep{SmartXIX} modifies the dynamics near but not at equilibrium, and its effect on probabilities is a potential experimental signature (Section~\ref{sec:future}).
\end{itemize}

The open questions concern the \emph{scope} of the conditional derivation: whether physical measurements realise the localised-probe-observer variant of Paper~XXIX (so that non-contextuality is a theorem rather than an extra axiom), whether the measurement-sampling hypothesis of Route~C holds in physical measurement contexts, whether the Laplace regime applies, and whether $\Mcl = \mathbb{CP}^{N-1}$ holds for the observers in question. Within those stipulations, the abstract assignment $P$ of Theorem~\ref{thm:born_unique} is uniquely fixed to the Born form.

\subsection{The equilibrium assumption}

The strongest route (C) depends on the system being \emph{at} the Paper~XIV equilibrium (Proposition~\ref{prop:born} hypothesis~(1) is stated as $\rho = \Pst$; a near-equilibrium version would require an additional error term not tracked here). This is analogous to the quantum-equilibrium hypothesis in Bohmian mechanics: both frameworks derive the Born distribution from a dynamical measure, and both require a justification for why the system is in the relevant equilibrium state. Paper~XIV establishes that the stochastic closure dynamics admits $\Pst \propto e^{-2\Fcal/\sigma^2}$ as a stationary distribution, with uniqueness contingent on irreducibility/regularity/confining-boundary conditions stated in its proof. Convergence of arbitrary initial distributions to $\Pst$ is not unconditionally established there; under additional ergodicity assumptions on the closure Markov process, one expects long-time convergence, but both the precise conditions and the relevant timescales are open questions.

%==================================================================
\section{Future Directions}\label{sec:future}
%==================================================================

\begin{enumerate}
    \item \textbf{Concrete constraint functional.} Constructing $\Ccal_\Ocal$ and the weight measure explicitly from the 600-cell geometry would move the framework from abstract construction to computable prediction.

    \item \textbf{Physical scope of the localised-probe-observer variant.} Paper~XXIX proves LMC within its localised-probe-observer class, and the present paper invokes this to derive non-contextuality (in restricted form). A remaining programme step is to argue that every physical measurement apparatus realises a localised probe observer in this sense --- that is, to show the probe decomposition, off-basin silence, and canonical-probe-assignment hypotheses hold for generic measurements in the VFD framework. If established, this would remove one observer-side input to Route~E (the scope condition), but not the separate completeness ($\Mcl(\Ocal) = \mathbb{CP}^{N-1}$) and basis-realisability assumptions needed to upgrade restricted non-contextuality to the global (U2).

    \item \textbf{Away-from-equilibrium corrections.} The nonlinear residual $\delta U_{\mathrm{rel}}$ (Paper~XIX) modifies the dynamics near but not at equilibrium. Its effect on outcome probabilities may produce measurable deviations from the Born rule --- a potential experimental signature.

    \item \textbf{Emergence of observer structures.} Showing how observer regions satisfying the Paper~XXIX definition arise naturally from the closure dynamics would close the last gap between geometry and measurement.

    \item \textbf{Decoherence integration.} The constraint framework should be compatible with decoherence as a mechanism for producing well-separated basins. Making this explicit would strengthen the assumptions of Proposition~\ref{prop:born}.
\end{enumerate}

%==================================================================
\section{Conclusion}\label{sec:conclusion}
%==================================================================

In this framework, probability is not a primitive stochastic ingredient added to measurement. It is the induced measure over the set of observer-compatible admissible configurations, with weights determined by the geometry and dynamics of constraint satisfaction.

The programme chain is:
\begin{enumerate}
    \item \textbf{Structure}: closure geometry on the 600-cell.
    \item \textbf{Dynamics}: stochastic closure process with stationary distribution $\Pst \propto e^{-2\Fcal/\sigma^2}$ (Paper~XIV).
    \item \textbf{Quantum recovery}: Schr\"odinger equation as equilibrium tangent limit, with $|\Psi|^2 = \Pst$ in the equilibrium regime (Papers~XVIII--XXI).
    \item \textbf{Observer}: constraint substructure determining admissible configurations (Paper~XXIX).
    \item \textbf{Measurement}: reduction to observer-admissible sectors (Paper~XXIX).
    \item \textbf{Probability (conditional existence, Route~C)}: in the equilibrium closure regime, with well-separated sectors, sector-supported decomposition, the sampling hypothesis (measurement frequencies equal integrated stationary sector masses), and state-completeness (the Nelson density is supported in the admissible set), the stationary sector measure gives $P_i = |c_i|^2$ up to overlap corrections.
    \item \textbf{Probability (conditional uniqueness, Route~E)}: Route~E's contribution is geometric packaging and motivation. A $U(N)$-invariant K\"ahler structure on $\mathbb{CP}^{N-1}$ is imported from Paper~XXI (as standard geometric QM on Paper~XXI's assumed invariant subspace) and identified with Fubini--Study up to scale by Petz's theorem; in the Laplace regime, the closure Hessian locally realises a configuration-space Fisher metric, which could supply a candidate local scale conditional on an unproved bridge theorem (Proposition~\ref{prop:fs_assembly}, Remark~\ref{rem:global_gap}). A restricted, state-specific form of non-contextuality is derived from basin locality of the closure flow together with a locality condition on observers (LMC); upgrading this to the global (U2) hypothesis of Theorem~\ref{thm:born_unique} requires \emph{both} universal completeness $\Mcl(\Ocal) = \mathbb{CP}^{N-1}$ \emph{and} the assumption that every orthonormal basis is realised by a complete localised-probe observer. The uniqueness proof itself is Gleason's under (U1)--(U4) and universal completeness for $N \geq 3$; the Hessian/Petz machinery does not enter the proof. The $N = 2$ case is not covered here (Remark~\ref{rem:N_two}).
\end{enumerate}

The result is a conditional recovery of the Born rule. Two independent arguments converge in the equilibrium tangent regime: Route~C realises the Born distribution as the stationary sector measure of the stochastic closure process (under Proposition~\ref{prop:born}'s five hypotheses, including the sampling postulate and the state-completeness condition), and Route~E shows that --- given the imported $U(N)$-invariant K\"ahler structure of Paper~XXI, the Laplace approximation, $\Mcl = \mathbb{CP}^{N-1}$, and axioms (U1)--(U4) of Theorem~\ref{thm:born_unique} --- no other sector-probability assignment satisfying those hypotheses is admissible. Non-contextuality of probability is adopted as the global hypothesis (U2); a restricted, state-specific form is derived within the complete localised-probe-observer variant of Paper~XXIX from basin locality and (LMC), and upgrades to the global (U2) only in combination with universal completeness \emph{and} the assumption that every orthonormal basis is realised by a complete localised-probe observer.

The assumptions under which this operates are stated explicitly: equilibrium, well-separated sectors, sector-supported decomposition, the Route~C sampling hypothesis, state-completeness, Laplace regime, measurement realised by an observer in Paper~XXIX's localised-probe-observer variant (within which LMC is a theorem), universal completeness $\Mcl(\Ocal) = \mathbb{CP}^{N-1}$, the basis-realisability assumption (every orthonormal basis is realised by a complete localised-probe observer) required to upgrade restricted (U2) to global (U2), and dimension $N \geq 3$ (with $N = 2$ not covered here; see Remark~\ref{rem:N_two}). The remaining observer-side inputs are: (a) the localised-probe-observer scope condition on physical measurements, (b) universal completeness, and (c) basis-realisability. The remaining open questions concern the \emph{scope} of these assumptions (when equilibrium obtains, how decoherence produces the sector structure, when the Laplace regime applies, whether every physical measurement is realised by a complete localised-probe observer for every orthonormal basis, and whether universal completeness holds), not the form of the weighting law within that scope.

Within the complete-observer equilibrium regime treated in Route~C, and under the stationary-sampling and sector-separation hypotheses, standard quantum predictions are recovered. What changes is the explanatory status of the Born rule: from primitive postulate to conditional consequence of stationary-sector weighting plus a Gleason uniqueness argument on imported geometric structure. Given Route~C's equilibrium, sector-separation, sampling, and state-completeness hypotheses plus Route~E's (U1)--(U4) plus universal completeness plus basis-realisability, the quadratic form is uniquely selected among smooth, non-contextual, $U(N)$-covariant probability assignments.

The framework also supplies an \emph{observer-dependent conditionalisation rule within the VFD observer formalism}: for \emph{partial} observers, whose admissible set omits some of the state's Nelson support, the framework assigns context-dependent conditional probabilities. Under the sector-supported decomposition regime of Route~C, this takes the explicit form $|c_i|^2 / \sum_{j \in \mathrm{adm}}|c_j|^2$, with the mass $1 - \sum_{j \in \mathrm{adm}}|c_j|^2$ interpreted as the probability of non-detection by that observer (Remarks~\ref{rem:partial}, \ref{rem:vs_standard_qm}). Outside the sector-supported regime, the conditional probability is still the normalised admissible-set integral of $|\Psi|^2$ (equation~\eqref{eq:prob_via_integration}), but the explicit sum-of-amplitudes form requires sector decomposition. Standard QM either hides this inside the POVM formalism or adds a no-click outcome by hand; in the present VFD observer formalism the same quantity is assigned by admissible-set normalisation. For complete observers the two prescriptions agree exactly; for partial observers they differ at the formalism level. Whether that difference has a physical realisation and genuine operational distinctness remains open, pending a measurement model that specifies when and how admissible-set normalisation is realised by actual detectors.

Outside these stated regimes, the question is open and the paper identifies where further work is needed. The conditional/partial-observer structure is not an artefact of the conditional derivation; it is a built-in feature of the present VFD observer formalism, whose physical scope remains open. Other interpretations recover partial-measurement probabilities only through extra machinery.

%==================================================================
\bibliographystyle{plainnat}
\bibliography{references}

\end{document}
